\documentclass[
  11pt, % The default document font size, options: 10pt, 11pt, 12pt
  %codirector, % Uncomment to add a codirector to the title page
]{charter}

% The title of the report, used on the cover page and available anywhere in the document via \ttitle
\titulo{Biometric identity verification for outsourced-staff management platforms}

% Name of the graduate program, used on the cover page and available anywhere in the document via \degreename
%\posgrado{Carrera de Especialización en Sistemas Embebidos}
%\posgrado{Carrera de Especialización en Internet de las Cosas}
\posgrado{Carrera de Especialización en Inteligencia Artificial}
%\posgrado{Maestría en Sistemas Embebidos}
%\posgrado{Maestría en Internet de las cosas}

% Your name, available anywhere in the document via \authorname
% IMPORTANT: do not omit degrees or accents in names; full legal names (as they appear on your ID) are recommended
\autor{Denis Zhbankov}

% Name of the director and co-director, available anywhere in the document via \supname, \cosupname, \pertesupname, and \pertecosupname
\director{Ksenija Blažević}
\pertenenciaDirector{Head of Machine Learning \& AI}
\codirector{} % uncomment the codirector option in the documentclass parameters for this to appear on the cover
\pertenenciaCoDirector{FIUBA}

% Name of the client who will approve the project results, available via \clientename and \empclientename
\cliente{Alina Morozovskaia}
\empresaCliente{Leroy Merlin}

\fechaINICIO{June 23, 2026}    %Fecha de inicio de la cursada de GdP \fechaInicioName
\fechaFINALPlan{August 18, 2026}   %Fecha de final de cursada de GdP
\fechaFINALTrabajo{April 23, 2027}  % Fecha de defensa pública del trabajo final

\begin{document}

\maketitle
\thispagestyle{empty}
\pagebreak

\thispagestyle{empty}
{\setlength{\parskip}{0pt}
  \tableofcontents{}
}
\pagebreak

\section*{Change log}
\label{sec:registro}

\begin{table}[ht]
  \label{tab:registro}
  \centering
  \begin{tabularx}{\linewidth}{@{}|c|X|c|@{}}
    \hline
    \rowcolor[HTML]{C0C0C0}
    Revision & \multicolumn{1}{c|}{\cellcolor[HTML]{C0C0C0}Details of changes made} & Date               \\ \hline
    0        & Document created                                                          & \fechaInicioName    \\ \hline
    1        & Completed through point 5 inclusive                                          & July 3, 2026  \\ \hline
    2        & Points 6 and 7 completed                                                   & July 10, 2026 \\ \hline
    3        & Everything completed                                                                & July 28, 2026 \\ \hline
  \end{tabularx}
\end{table}

\pagebreak

\section*{Project charter}
\label{sec:acta}

\begin{flushright}
  Buenos Aires, \fechaInicioName
\end{flushright}

\vspace{2cm}

By this document it is agreed with \authorname\hspace{1px} that their final graduate project for the \degreename\hspace{1px} will be titled ``\ttitle'' and will consist of the design, domain-adaptation fine-tuning, and evaluation of a biometric verification module that, through facial recognition, allows the identity of workers to be checked at the start of each work shift and against the photograph on the identity document presented, thereby helping to verify identity correspondence within the external workforce. The work is carried out in an academic setting (evaluation with public \textit{datasets}, fine-tuning of pretrained models, and integration via a CLI \textit{harness} and an interactive \textit{demo}) and has a preliminary estimated budget of 600 hours.

The project budget is broken down into two clearly differentiated components:

\begin{itemize}
  \item Engineering cost (included in the charter): USD 36,000, corresponding to 600 professional hours at a reference market rate of USD 60/h for AI consulting projects of this scale. This amount is listed in the charter for reference purposes only and does not imply an actual charge to the client under the academic agreement.
  \item Contingent compute operating costs (not included in the charter): up to USD 1,000 over the course of the project, for a low-cost compute subscription (e.g., Colab Pro, USD 10/month) and academic or cloud \textit{credits} as a contingency. These are reported separately from the engineering cost, detailed in the budget table below, and justified in section 13 as part of the mitigation for Risk 1.
\end{itemize}

\begin{table}[ht]
  \centering
  \caption{Itemized project budget.}
  \label{tab:presupuesto}
  \begin{tabularx}{\linewidth}{@{}|X|c|c|r|@{}}
    \hline
    \rowcolor[HTML]{C0C0C0}
    Component & Quantity & Rate & Subtotal \\
    \hline
    \multicolumn{4}{|l|}{\cellcolor[HTML]{EEEEEE}\textbf{Engineering cost (included in the charter)}} \\
    \hline
    Work hours (development/engineering) & 600 h & USD 60/h & USD 36,000 \\
    \hline
    \multicolumn{3}{|r|}{\textbf{Engineering cost subtotal}} & \textbf{USD 36,000} \\
    \hline
    \multicolumn{4}{|l|}{\cellcolor[HTML]{EEEEEE}\textbf{Contingent operating costs (not included in the charter)}} \\
    \hline
    Compute subscription (Colab Pro, primary path) and academic or cloud \textit{credits} (contingency) & 1 & --- & USD 1,000 \\
    \hline
    \multicolumn{3}{|r|}{\textbf{Contingent operating costs subtotal}} & \textbf{USD 1,000} \\
    \hline
  \end{tabularx}
\end{table}

The work start date is \fechaInicioName\hspace{1px} and the public presentation date is \fechaFinalName.

The initial plan is attached to this charter.

\vfill

% This part is built automatically from the information loaded in the document preamble and should not be modified
\begin{table}[ht]
  \centering
  \begin{tabular}{ccc}
    \begin{tabular}[c]{@{}c@{}}Dr. Ing. Ariel Lutenberg \\ FIUBA Graduate Program Director
    \end{tabular} & \hspace{2cm} &
    \begin{tabular}[c]{@{}c@{}}\clientename \\ \textit{Product owner}, external company (NDA)
    \end{tabular} \vspace{2.5cm} \\
    \multicolumn{3}{c}{
      \begin{tabular}[c]{@{}c@{}} \supname \\ Final project director
    \end{tabular}} \vspace{2.5cm} \\
  \end{tabular}
\end{table}

\section{1. Technical-conceptual description of the project to be carried out}
\label{sec:descripcion}

This work is carried out within \empclientename, an international home-improvement company, which faces a recurring problem in the management of its outsourced operational staff (warehouse workers, cleaning staff, mechanics, laborers): the lack of certainty that the person present is the same one listed in the documentation or in the employment record. In practice, the operational risk is concentrated in two identity-correspondence situations: (1) during hiring, that the person presenting themselves to the recruiter does not match the photograph on the identity document shown; and (2) in daily operations, that the person attending a shift is not the same person who was registered as the worker. Outside the scope of this work is determining whether the identity document is authentic or forged, as well as verifying that the captured image corresponds to a live person rather than a reproduction.

Regarding the state of the art, facial recognition is a mature field with reference architectures such as ArcFace and AdaFace for generating facial \textit{embeddings}, and detection models such as MTCNN and RetinaFace. Additionally, TransFace (ICCV 2023) represents a recent \textit{transformer}-based extension that reports incremental improvements over ArcFace on recent \textit{benchmarks}; it is included as an optional evaluation model if the schedule allows (see section 4.1). The central question motivating this work is no longer which pretrained model is abstractly better (that question has a known answer in the literature), but rather what effect domain-adaptation \textit{fine-tuning} (which may be positive, neutral, or negative) has on the verification performance of a facial recognition model, when applied to a subset of a public \textit{dataset} augmented with synthetic degradations that emulate adverse operational conditions (blur, compression, low resolution) and age-diversity-oriented sampling, and what the associated computational cost \textit{trade-off} is. This question is relevant because generic solutions do not account for the particular operational conditions of the operational work context: heterogeneous capture-device quality, uncontrolled lighting, aging between the document photo and the current image, and variable quality of digitized documents.

The proposed solution consists of a biometric verification module based on domain-adaptation \textit{fine-tuning} of modern pretrained models (ArcFace and AdaFace as the main models, TransFace as optional) on a subset of a large-scale public \textit{dataset} (Glint360K or WebFace4M) augmented with synthetic degradations (blur, compression, low resolution) and with age-diversity-oriented sampling, in order to bring the training subset closer to the adverse operational conditions of the target domain, together with its systematic evaluation on three standard \textit{benchmarks} with complementary conditions: LFW (classic, saturated \textit{benchmark}, expected to exceed 99\% on both models), AgeDB (\textit{cross-age} verification, which directly addresses the stated aging problem), and RFW (\textit{Racial Faces in the Wild}, which measures demographic bias and ties into the governance section). The project is carried out in an academic context without direct access to the company's operational data; the integration is materialized through a CLI \textit{harness} that runs the full \textit{pipeline} on the \textit{benchmarks} and produces reproducible metrics, and an interactive \textit{demo} built with Gradio for the presentation. To close the loop between training and evaluation under adverse conditions, the same synthetic degradations applied to the \textit{fine-tuning} subset are also applied to build degraded variants of LFW and AgeDB, which emulate the capture-quality heterogeneity of the operational domain and allow comparable measurement of performance under controlled adverse conditions, for both the pretrained and the fine-tuned model.

The core technical \textit{pipeline} consists of: (1) facial detection using models such as MTCNN or RetinaFace; (2) \textit{landmark}-based alignment; (3) generation of facial \textit{embeddings} using ArcFace or AdaFace, in either their pretrained configuration or after \textit{fine-tuning} on the selected subset; (4) similarity computation via cosine distance between the \textit{embedding} of the person present and the reference \textit{embedding}; (5) a thresholded decision that balances the false-positive and false-negative rates according to the security requirements of the scenario; and (6) a minimal \textit{liveness} module based on a pretrained model, operating as an extensible \textit{baseline} (not as a production system). An interactive \textit{demo} is also included to showcase the registration and verification flows for the presentation.

Figure \ref{fig:diagBloques} presents the system's block diagram, distinguishing the prototype's two functional flows. The registration flow (document--face verification) starts from two inputs: the live capture from the device camera and the digitized identity document. Both images go through the stages of face detection with quality validation (a single face, sufficient sharpness and lighting), \textit{landmark}-based alignment, and generation of the facial \textit{embedding} using the selected model (ArcFace or AdaFace). The resulting pair of \textit{embeddings} is compared via cosine distance against a configurable threshold, while, in parallel, an auxiliary \textit{liveness} classifier evaluates the live capture and contributes its \textit{score} to the decision. If the decision is affirmative, the validated \textit{embedding} is stored as the reference in the database; otherwise, a new capture is requested or the case is referred to manual review. The shift-verification flow starts only from the live capture, goes through the same processing stages, and compares the resulting \textit{embedding} against the stored reference, with \textit{liveness} verification as an additional condition for the decision. On rejection, the system allows up to three retries, logs the incident for auditing, and refers the case to manual review.

\begin{figure}[htpb]
  \centering
  \includegraphics[height=0.985\textheight,keepaspectratio]{./Figuras/diagBloques.png}
  \caption{System block diagram.}
  \label{fig:diagBloques}
\end{figure}

Since this is a collaborative project with a company within the academic framework of the graduate program, the project does not include funding from additional public or private programs. The relationship with the client is governed by a non-disclosure agreement (NDA) between the \textit{product owner} and the external company they belong to; this NDA covers commercial information and internal data of the client and of the adopting company, and does not restrict the dissemination of the work's academic results: the methodology, the prototype code, the public \textit{datasets} used, the evaluation metrics, and the experimental analysis are publishable as part of the final project. The documentation avoids including confidential data belonging to the client or to \empclientename.

The differential value of this project rests on three aspects. First, it addresses a concrete application domain (workplace identity verification in operational outsourcing) with particular operational conditions that generic facial-recognition solutions do not account for. Second, it proposes a systematic experimental analysis of the impact of domain-adaptation \textit{fine-tuning} (using a subset augmented with the same synthetic degradations applied to the evaluation \textit{benchmarks}) on performance and computational cost, an approach rarely reported in the recent comparative literature and one that, unlike fine-tuning on an untransformed subset, allows a genuinely answerable experimental question to be formulated. Third, it integrates fairness and demographic-robustness considerations from the design stage (evaluation on RFW) and lays the groundwork for an extensible \textit{liveness} module in a later iteration (for example, at a master's-degree stage), and it relies exclusively on public \textit{datasets} subject to verification of their licenses and terms of use, in order to ensure compliance with personal-data-protection regulations within the scope of the final project.

\section{2. Identification and analysis of stakeholders}
\label{sec:interesados}

In a project it is a priority to know and manage the stakeholders, their expectations, and their influence. This section focuses on clearly detailing the role each person involved must play throughout the entire process.

\begin{table}[ht]
  \begin{tabularx}{\linewidth}{@{}|l|X|X|l|@{}}
    \hline
    \rowcolor[HTML]{C0C0C0}
    Role           & Full name & Organization    & Position        \\ \hline
    Client       & \clientename      & External company (NDA) & \textit{Product owner} \\ \hline
    Advisor   & \supname          & \empclientename & \pertesupname                     \\ \hline
    Responsible party   & \authorname       & FIUBA           & Student                            \\ \hline
    End user & ---               & \empclientename operational profiles & HR specialist, shift supervisor \\ \hline
  \end{tabularx}
\end{table}

The project is framed within \empclientename\hspace{1px} as the adopting company and owner of the problem, institutionally represented by the advisor. The role of client (guiding business value and approving deliverables) is held by \clientename, \textit{product owner} of an external company under an NDA relationship, with no corporate ties to \empclientename. This configuration arises because the relationship with \empclientename\hspace{1px} within the framework of this final project is channeled through the advisor as the institutional point of contact, while the client operates from a separate organization with its own confidentiality agreement, which is documented explicitly in order to delimit the scope of the validation carried out with public \textit{datasets}.

Main characteristics of the stakeholders:

\begin{itemize}
  \item Client (\clientename)
    \begin{itemize}
      \item Role: acts as \textit{product owner} under NDA.
      \item Expectations: that the prototype demonstrates the verification flows, that the experimental analysis is rigorous, and that the results provide useful technical evidence for a possible future deployment decision.
      \item Influence: high --- sets \textit{backlog} priorities and approves closure.
      \item Information received: periodic progress updates, evaluation results, technical documentation.
      \item Deliverables reviewed: prioritized \textit{backlog}, evaluation \textit{harness} and results, documentation, final project deliverable.
    \end{itemize}
  \item Advisor (\supname)
    \begin{itemize}
      \item Role: \textit{machine learning} expert. Oversees the definition of technical requirements, architecture selection, \textit{fine-tuning}, and evaluation of the biometric \textit{pipeline}.
      \item Expectations: that technical decisions are well founded and that the experimental evaluation is rigorous and reproducible.
      \item Influence: high --- oversees the technical direction and validates the choice of models and protocols.
      \item Information received: requirements, experiment results, performance metrics.
      \item Deliverables reviewed: architecture selection, \textit{fine-tuning} configuration, comparative report, evaluation \textit{harness}.
    \end{itemize}
  \item Responsible party (\authorname)
    \begin{itemize}
      \item Role: in charge of planning, implementing, tuning, evaluating, and documenting the prototype through the final defense. In the academic context of the work, the role of prototype administrator that appears in the user stories (HU7, HU14, and HU15) is carried out by the responsible party themself.
      \item Expectations: to complete the final project with quality, within the estimated time frame and hours.
      \item Influence: medium --- executes and documents, but depends on decisions made by the client and the advisor.
      \item Information received: meetings with the client and the advisor, feedback on deliverables.
      \item Deliverables reviewed: all the code, the technical documentation, and the final project report.
    \end{itemize}
  \item End user (HR specialist, shift supervisor)
    \begin{itemize}
      \item Role: the reference operational roles for the prototype are the HR specialist (registration flow and document--face verification) and the shift supervisor (start-of-shift verification flow). This work does not deploy the system on \empclientename's real operations; these roles guide the design of the \textit{demo} and of the user stories.
      \item Expectations: that the system be intuitive, fast, and reliable for the registration and verification flows.
      \item Influence: low --- reference users of the prototype; they do not take part in project decisions.
      \item Information received: interactive \textit{demo}, user manuals.
      \item Deliverables reviewed: \textit{demo} and usage experience during testing with public \textit{datasets}.
    \end{itemize}
\end{itemize}

\section{3. Project purpose}
\label{sec:proposito}

The purpose of this project is to design, tune, evaluate, and document, within an academic framework, a biometric verification module based on facial recognition, to help verify the identity correspondence of outsourced staff at the start of the shift and against the photograph on the document presented. The central experimental question of the work is what effect domain-adaptation \textit{fine-tuning} (which may be positive, neutral, or negative) has on the verification performance of a modern model (ArcFace, AdaFace), when that tuning is applied to a subset of a public \textit{dataset} augmented with synthetic degradations that emulate adverse operational conditions (blur, compression, low resolution) and with age-diversity-oriented sampling, and what the \textit{trade-off} with the associated computational cost is.

\paragraph{Primary metric and \textit{benchmark}.}
The primary experimental-success metric is the \textit{Equal Error Rate} (EER) on AgeDB, and the primary \textit{benchmark} is AgeDB in its standard configuration, since it directly addresses the aging problem between the document photo and the live capture stated in the project description. The secondary metrics are FAR, FRR, AUC, TAR at fixed FAR, mean inference time, and memory usage, evaluated on LFW, RFW, and the degraded variants of LFW and AgeDB.

The MVP components and the optional ones, the included and excluded scope, and the details of the registration and verification flows are detailed in sections 4 and 4.1.

\section{4. Project scope}
\label{sec:alcance}

The project includes:

\begin{itemize}
  \item implementation of a biometric verification \textit{pipeline} based on modern pretrained models (ArcFace, AdaFace): face detection, alignment, \textit{embedding} generation, and similarity comparison;
  \item verification of correspondence between the live face and the photograph on the document in the registration flow, and facial verification against the stored reference at the start of the shift;
  \item biometric comparison module with configurable thresholding to balance security and usability;
  \item domain-adaptation \textit{fine-tuning} of at least one of the models (ArcFace or AdaFace) on a selected subset of a large-scale public \textit{dataset} (Glint360K or WebFace4M), augmented with synthetic degradations (blur, compression, low resolution) and age-diversity-oriented sampling, in order to bring the training subset closer to the adverse operational conditions of the target domain, with an assessment of its impact relative to the pretrained version;
  \item comparative evaluation of the pretrained and fine-tuned models on at least three public \textit{benchmarks}: LFW (\textit{baseline}), AgeDB (\textit{cross-age}), and RFW (demographic bias), as well as on synthetically degraded variants of LFW and AgeDB (blur, compression, low resolution), with EER, FAR, FRR, ROC curve, AUC, TAR at fixed FAR (e.g., TAR at FAR = $10^{-3}$), and mean inference time metrics. The main models are ArcFace and AdaFace; TransFace (ICCV 2023) may be included as a third, optional model if the schedule allows, in order to broaden the state-of-the-art coverage (see section 4.1);
  \item \textit{liveness} module (mandatory MVP component) based on a pretrained model (e.g., MiniFASNet or FeatherNet) evaluated on a public \textit{spoofing} \textit{dataset} (OULU-NPU or CASIA-FASD), operating as an extensible \textit{baseline} whose robust development is left as future work;
  \item \textit{end-to-end} evaluation \textit{harness}, runnable as a CLI or a parameterizable \textit{notebook}, that runs the full \textit{pipeline} on a \textit{benchmark} and produces reproducible metrics and figures;
  \item interactive \textit{demo} with Gradio or Streamlit to present the prototype;
  \item documentation including: a progress report, a technical report, a prototype user manual, \textit{harness} documentation, and a repository with instructions for reproducing the experiments.
\end{itemize}

This project does not include:

\begin{itemize}
  \item validation of the authenticity, integrity, or validity of identity documents (detection of forged documents, tampering with physical media, security OCR, official verification \textit{chips} or codes);
  \item production-grade robust \textit{liveness detection}: the included module operates as an academic \textit{baseline} and not as a certified anti-\textit{spoofing} system. Its robust development is left as a future line of work;
  \item training facial-recognition models from scratch: the project starts from pretrained weights and, where applicable, performs limited \textit{fine-tuning} on a subset of the \textit{dataset};
  \item collection of biometric images of workers or of the company's internal data. Evaluation will be carried out exclusively with public or synthetic \textit{datasets} authorized for academic use;
  \item development of \empclientename's\hspace{1px} staff-management platform or of its existing functionality;
  \item production deployment, integration with the company's real systems, or post-delivery maintenance;
  \item development of capture \textit{hardware} or camera procurement;
  \item external certification of the \textit{liveness} module or a formal security audit;
  \item certification of compliance with HIPAA, ISO, or other regulatory frameworks, although an analysis of the regulatory compliance applicable to the processing of biometric and personal data is included (section 12);
  \item implementation of access-control mechanisms for the \textit{embedding} store, deletion or updating of biometric records, or informed-consent infrastructure: these measures are documented in section 12 (data governance) as requirements for a possible future deployment and are not part of the prototype.
\end{itemize}

\subsection{MVP components and optional components}

The mandatory MVP components are those whose delivery is indispensable to answer the work's experimental question. The optional components broaden the state-of-the-art coverage or reinforce some aspect of the prototype and are carried out only if the schedule allows after completing the MVP. The separation is kept traceable throughout the \textit{product backlog} (section 6), the WBS (section 9), the sprint plan (section 10), and the compute budget (project charter, table~\ref{tab:presupuesto}).

\paragraph{Mandatory components (MVP).}
\begin{itemize}
  \item Biometric verification \textit{pipeline} with pretrained ArcFace and AdaFace (detection, alignment, \textit{embedding}, similarity comparison).
  \item Domain-adaptation \textit{fine-tuning} of ArcFace on a subset of Glint360K or WebFace4M, evaluated on LFW, AgeDB, and RFW and on the degraded variants of LFW and AgeDB.
  \item \textit{Liveness} module integrated into the \textit{pipeline} as an academic \textit{baseline}, evaluated on OULU-NPU or CASIA-FASD and reported together with its limitations.
  \item Runnable, reproducible CLI \textit{harness}.
  \item Interactive \textit{demo} with Gradio or Streamlit.
  \item Complete technical documentation (progress report, thesis report, user manual, \textit{harness} documentation, repository with reproduction instructions).
\end{itemize}

\paragraph{Optional components (not part of the MVP).}
\begin{itemize}
  \item Evaluation of TransFace (ICCV 2023) as a third pretrained model, under identical conditions to ArcFace and AdaFace (8 h in the WBS).
  \item Domain-adaptation \textit{fine-tuning} of AdaFace as a second tuning pass, using the same MVP methodology (12 h in the WBS).
  \item Any extension of the \textit{liveness} module beyond the \textit{baseline} (alternative architectures, tuning on domain \textit{spoofing} data): documented as future work.
\end{itemize}

The 20 h of optional tasks (TransFace as a third model + a second \textit{fine-tuning} pass) are reflected in the WBS (section 9) and excluded from the sprint plan, so that their execution remains strictly subordinate to time availability after the MVP is completed.

\section{5. Project assumptions}
\label{sec:supuestos}

For the development of this project, it is assumed that:

\begin{itemize}
  \item Public \textit{datasets} of sufficient quantity and quality will be accessible for the evaluation and \textit{fine-tuning}, within the project's limitations. In particular, Glint360K (or WebFace4M as an alternative) will be available for academic download under a compatible license.
  \item The license terms of each public \textit{dataset} used will be verified before use and documented in the data-governance section.
  \item Sufficient computational resources are available for experimentation, inference, and limited \textit{fine-tuning} of the models, within the free-tier usage limits of common compute platforms (Google Colab, Kaggle, Lightning AI) or through a low-cost subscription. Purchasing decisions, the auxiliary budget, and the response to a potential shortfall in compute are addressed in section 13 (risk management) and reflected in the project charter (table~\ref{tab:presupuesto}); they are not part of the project's assumptions.
  \item The time needed to complete the estimated 600 hours of the project is available within the established time frame.
  \item Macroeconomic and regulatory conditions will not undergo significant changes that affect the project's execution or budget.
  \item The technical skills needed in Python, PyTorch, and computer vision to implement the solution are available. A learning curve at the start of the project is expected, since the responsible party is new to facial recognition; this is offset by choosing well-documented models and tools (InsightFace, \textit{adaface-pytorch}, Gradio).
  \item The \textit{fine-tuning} is carried out on a representative subset of a public \textit{dataset}, not on the company's real data, so the experimental question is scoped to improvement under synthetic domain conditions.
  \item Scope changes will be exceptional and will be managed by agreement with the advisor and the client.
\end{itemize}

\section{6. \textit{Product backlog}}
\label{sec:backlog}

The \textit{product backlog} is organized into four epics aligned with the identity-correspondence problem, the academic scope of the prototype, and the expected experimental results. The stories describe functionality of the demonstration module; they do not imply operational deployment at \empclientename.

\subsection{Weighting criterion: \textit{story points}}

Each user story is assessed along three dimensions:

\begin{itemize}
  \item Difficulty: estimated amount of work (scale of 1 to 8).
  \item Complexity: required technical level (scale of 1 to 8).
  \item Uncertainty: associated risk or novelty (scale of 1 to 8).
\end{itemize}

The sum of the three dimensions is rounded up to the nearest number in the Fibonacci series: 1, 2, 3, 5, 8, 13, 21, 34\dots \textit{Story points} are relative and do not directly correspond to hours of work.

\subsection{Epic 1: Document--face registration and verification}

HU1: Facial capture of the candidate
\begin{itemize}
  \item As an HR specialist, I want to capture a candidate's face using the device camera, in order to generate and store their biometric representation as a future reference.
  \item Priority: high
  \item Difficulty: 5, complexity: 5, uncertainty: 4
  \item \textit{Story points}: 21
\end{itemize}

HU2: Face extraction from the identity document
\begin{itemize}
  \item As an HR specialist, I want the system to extract the face from the photograph on the identity document (national ID, driver's license, resident card) so it can be used in the biometric cross-check at registration.
  \item Priority: high
  \item Difficulty: 5, complexity: 5, uncertainty: 5
  \item \textit{Story points}: 21
\end{itemize}

HU3: Document--face verification during hiring
\begin{itemize}
  \item As an HR specialist, I want the system to compare the live-captured face with the face extracted from the document and issue an acceptance or rejection decision.
  \item Priority: high
  \item Difficulty: 5, complexity: 6, uncertainty: 5
  \item \textit{Story points}: 21
\end{itemize}

HU4: Handling invalid captures
\begin{itemize}
  \item As an HR specialist, I want the system to detect and flag anomalous situations during capture (no face detected, multiple faces, blurry image, insufficient lighting, partially occluded face) and request a new capture.
  \item Priority: high
  \item Difficulty: 5, complexity: 5, uncertainty: 3
  \item \textit{Story points}: 13
\end{itemize}

\subsection{Epic 2: Identity verification at the start of the shift}

HU5: Facial verification at the start of the shift
\begin{itemize}
  \item As a shift supervisor, I want the system to verify, at the start of each shift via facial recognition against the stored reference, that the person present matches the one registered.
  \item Priority: high
  \item Difficulty: 8, complexity: 5, uncertainty: 5
  \item \textit{Story points}: 21
\end{itemize}

HU6: Retry and rejection handling
\begin{itemize}
  \item As a shift supervisor, I want the system to allow up to three retries on rejection or an invalid capture and, if it persists, to refer the case to manual review and log the incident.
  \item Priority: medium
  \item Difficulty: 3, complexity: 3, uncertainty: 3
  \item \textit{Story points}: 13
\end{itemize}

\subsection{Epic 3: Experimental evaluation and \textit{fine-tuning}}

This epic concentrates the experimental core of the work. Its goal is to produce reproducible results on standard \textit{benchmarks} that answer the project's central question.

HU7: Acceptance-threshold configuration
\begin{itemize}
  \item As the prototype administrator, I want to select and configure a previously evaluated cosine-distance threshold, in order to balance security (FAR) and usability (FRR) according to the reference scenario.
  \item Priority: high
  \item Difficulty: 5, complexity: 5, uncertainty: 3
  \item \textit{Story points}: 13
\end{itemize}

HU8: Comparative evaluation of pretrained ArcFace and AdaFace
\begin{itemize}
  \item As the technical advisor, I want a comparative analysis of ArcFace and AdaFace in their pretrained configuration on the LFW, AgeDB, and RFW \textit{benchmarks} and on the degraded variants of LFW and AgeDB, with EER, FAR, FRR, AUC, TAR at fixed FAR, ROC curve, mean inference time, and memory usage metrics, in order to establish the experimental \textit{baseline}.
  \item Priority: high
  \item Difficulty: 8, complexity: 8, uncertainty: 4
  \item \textit{Story points}: 21
\end{itemize}

HU9: Evaluation of pretrained TransFace (optional, not part of the MVP)
\begin{itemize}
  \item As the technical advisor, I want that, if the schedule allows, TransFace (ICCV 2023) be evaluated under identical conditions to ArcFace and AdaFace and reported in an additional row of the comparison tables, with an explicit note that it results from an optional iteration and does not affect the work's main findings.
  \item Priority: low
  \item Difficulty: 4, complexity: 4, uncertainty: 4
  \item \textit{Story points}: 13
\end{itemize}

HU10: Preparation of the domain-adaptation \textit{fine-tuning} subset
\begin{itemize}
  \item As the technical advisor, I want a subset of Glint360K or WebFace4M (5,000--10,000 identities, 100,000--300,000 images) to be selected and documented, using age-diversity-oriented sampling, and for \textit{data augmentation} with synthetic degradations (blur, compression, low resolution) to be applied along with a documented \textit{train}/\textit{val} split, in order to bring the training subset closer to the adverse operational conditions of the target domain.
  \item Priority: high
  \item Difficulty: 5, complexity: 4, uncertainty: 4
  \item \textit{Story points}: 13
\end{itemize}

HU11: Execution of ArcFace's domain-adaptation \textit{fine-tuning}
\begin{itemize}
  \item As the technical advisor, I want ArcFace to be tuned via domain-adaptation \textit{fine-tuning} on the subset prepared in HU10, with a documented layer-\textit{freeze} scheme, \textit{learning rate}, number of epochs, and \textit{batch size}, and with intermediate \textit{checkpoints} preserved, in order to quantify the effect of the tuning (positive, neutral, or negative) on the \textit{cross-age} and degraded \textit{benchmarks}.
  \item Priority: high
  \item Difficulty: 8, complexity: 8, uncertainty: 5
  \item \textit{Story points}: 21
\end{itemize}

HU12: Execution of AdaFace's domain-adaptation \textit{fine-tuning} (optional, not part of the MVP)
\begin{itemize}
  \item As the technical advisor, I want that, if the schedule allows, AdaFace be tuned via domain-adaptation \textit{fine-tuning} on the same subset prepared in HU10, using the same methodology as HU11, in order to broaden the pretrained-vs-tuned comparison.
  \item Priority: low
  \item Difficulty: 4, complexity: 5, uncertainty: 4
  \item \textit{Story points}: 13
\end{itemize}

HU13: Post-\textit{fine-tuning} re-evaluation and comparative analysis
\begin{itemize}
  \item As the technical advisor, I want the models tuned in HU11 (and, if applicable, HU12) to be re-evaluated on LFW, AgeDB, RFW, and the degraded variants of LFW and AgeDB under identical conditions to HU8, and for tables, overlaid ROC curves, and a \textit{trade-off} analysis between accuracy and computational cost to be produced, in order to support the selection of the prototype's model.
  \item Priority: high
  \item Difficulty: 5, complexity: 6, uncertainty: 3
  \item \textit{Story points}: 21
\end{itemize}

HU14: \textit{Liveness} module (\textit{baseline})
\begin{itemize}
  \item As the prototype administrator, I want to integrate a pretrained \textit{liveness} classifier (e.g., MiniFASNet or FeatherNet) capable of discriminating between a genuine face and a presentation attack (printed photo, \textit{replay}), evaluated on a public \textit{spoofing} \textit{dataset} (OULU-NPU or CASIA-FASD), in order to have a documented, extensible \textit{baseline}.
  \item Priority: medium
  \item Difficulty: 3, complexity: 5, uncertainty: 4
  \item \textit{Story points}: 13
\end{itemize}

\subsection{Epic 4: Evaluation \textit{harness}, \textit{demo}, and delivery}

HU15: \textit{End-to-end} evaluation \textit{harness}
\begin{itemize}
  \item As the prototype administrator, I want a \textit{harness} runnable as a CLI or a parameterizable \textit{notebook} that runs the full \textit{pipeline} (detection, alignment, \textit{embedding}, comparison, decision) on a selected \textit{benchmark} and produces reproducible metrics and figures (ROC, EER, FAR, FRR, AUC, inference time), in order to allow independent verification of the results.
  \item Priority: high
  \item Difficulty: 5, complexity: 5, uncertainty: 2
  \item \textit{Story points}: 13
\end{itemize}

HU16: Interactive \textit{demo} and documentation
\begin{itemize}
  \item As an end user (HR or shift supervisor), I want an interactive \textit{demo} with Gradio or Streamlit that lets me run the registration and verification flows and view the outcome, along with usage documentation for the \textit{harness} and the \textit{demo}, in order to have clear material for the presentation and for any future replication.
  \item Priority: medium
  \item Difficulty: 3, complexity: 2, uncertainty: 2
  \item \textit{Story points}: 8
\end{itemize}

\section{7. User story acceptance criteria}
\label{sec:criteriosAceptacion}

Acceptance criteria are defined for each user story in the \textit{product backlog}, so that the conditions needed to consider the functionality complete can be validated. Each story has measurable, specific, and verifiable criteria. At least three criteria are included per story: one functional, one visual, and one technical (where applicable). In HU8, HU9, HU10, HU11, HU12, and HU13, the ``visual'' criterion is interpreted as a clear presentation of results (tables or charts in the documentation or in the \textit{demo}).

Each technical criterion specifies, where relevant, the test set, the experimental protocol, and the reference \textit{hardware}, in order to ensure reproducibility. The reference \textit{hardware} for all timing metrics is a 2019 16'' MacBook Pro (Intel, 16 GB RAM, AMD Radeon Pro 5300M 4 GB). \textit{Fine-tuning} and evaluation runs are preferably carried out in a CUDA-capable GPU environment (Google Colab, Kaggle, or equivalent), in order to keep training times reasonable within the project's schedule. Note that inference times measured on the reference \textit{hardware} (CPU/IGPU) are useful as reproducible bounds for the academic prototype and are not representative of a production deployment on GPU; the report documents the device used for each run, so that the figures are traceable and comparable.

\subsection{Epic 1: Document--face registration and verification}

Acceptance criteria for HU1: Facial capture of the candidate
\begin{enumerate}
  \item Functional: the system allows entering the candidate's data (name, ID number, position), capturing their facial image via the device camera, and storing the generated \textit{embedding}, associating it with the correct candidate.
  \item Visual: the interface shows a preview of the captured image with the face outlined before confirming registration.
  \item Technical: exactly one face is detected in the captured image with a confidence level above 80\%, the face is aligned, and the image has a minimum resolution of 160$\times$160 pixels and a sharpness measure (normalized Laplacian variance over the facial crop) above 0.05. The 0.05 threshold is a configurable parameter of the quality module, with an initial value of 0.05 and validation against the test set defined in HU4. The resulting \textit{embedding} is stored successfully and can be retrieved for later comparisons.
\end{enumerate}

Acceptance criteria for HU2: Face extraction from the identity document
\begin{enumerate}
  \item Functional: the system allows uploading an image of an identity document (national ID, driver's license, resident card) and automatically extracts the holder's photograph.
  \item Visual: a preview of the image extracted from the document is shown before confirming registration, so the user can validate its quality.
  \item Technical: on a set of at least 50 document images built from public \textit{datasets} or synthetic data (with variations in document type, capture angle, and lighting), the extraction detects and crops the holder's face with an \textit{IoU} above 70\% relative to the manual annotation in at least 90\% of cases. The resulting image is compatible with the \textit{embedding}-generation \textit{pipeline}.
\end{enumerate}

Acceptance criteria for HU3: Document--face verification during hiring
\begin{enumerate}
  \item Functional: the system compares the live face \textit{embedding} with that of the face extracted from the document and issues an acceptance or rejection result based on the configured threshold.
  \item Visual: the interface displays both images (or facial crops), the match result (accepted/rejected), and the similarity level.
  \item Technical: on pairs derived from LFW or controlled synthetic pairs, the match achieves an EER below 10\% and the total \textit{pipeline} time (detection + alignment + \textit{embedding} generation + comparison) is under 5 seconds on the reference \textit{hardware}.
  \item Error handling: when no face is present, multiple faces are detected, or the image does not meet the minimum quality, the system issues a descriptive message and requests a new capture without proceeding to the comparison.
\end{enumerate}

Acceptance criteria for HU4: Handling invalid captures
\begin{enumerate}
  \item Functional: the system detects the following anomalous situations and notifies the user with a descriptive message: no face detected, multiple faces detected, blurry image, insufficient lighting, partially occluded face.
  \item Visual: each error type is visually distinguished (e.g., with a specific message or icon), and the interface indicates the expected corrective action.
  \item Technical: image-quality detection (sharpness, lighting, occlusion) is performed before \textit{embedding} generation, avoiding processing of invalid inputs. On the test set with at least 20 images per error type, the correct anomaly-detection rate is above 85\%.
\end{enumerate}

\subsection{Epic 2: Identity verification at the start of the shift}

Acceptance criteria for HU5: Facial verification at the start of the shift
\begin{enumerate}
  \item Functional: the system detects the worker's face, generates its \textit{embedding}, compares it with the stored one, and issues an acceptance or rejection result.
  \item Visual: the \textit{demo} shows the camera with the detected face and, after verification, a green (accepted) or red (rejected) indicator along with the similarity level.
  \item Technical: the full \textit{pipeline} (detection, alignment, \textit{embedding} generation, comparison, decision) runs in under 5 seconds on the reference \textit{hardware}, following the verification protocol defined with the evaluation \textit{benchmark}.
  \item Error handling: when no face, multiple faces, or insufficient quality is present, the system requests a new capture without issuing a decision.
\end{enumerate}

Acceptance criteria for HU6: Retry and rejection handling
\begin{enumerate}
  \item Functional: on rejection or an invalid capture, the system allows up to three retries. If verification remains rejected after the retries, the \textit{demo} indicates that the case is referred to manual review and logs the incident with a \textit{timestamp} and attempt details.
  \item Visual: each retry shows a counter (attempt 1 of 3, etc.) and, after the third rejection, a referral message is shown.
  \item Technical: each attempt (rejection or invalid) is logged in the audit file with a \textit{timestamp}, the worker's identifier, the failure type, and the similarity level where applicable.
\end{enumerate}

\subsection{Epic 3: Experimental evaluation and \textit{fine-tuning}}

Acceptance criteria for HU7: Acceptance-threshold configuration
\begin{enumerate}
  \item Functional: the administrator can select among previously evaluated thresholds (results from HU8, HU11, and HU13) and apply them to subsequent verifications. The panel indicates the effect of each threshold (stricter: lower FAR, higher FRR; more permissive: higher FAR, lower FRR).
  \item Visual: the configuration panel shows the current threshold, a comparison table of evaluated thresholds with their associated metrics, and a control to save the selection.
  \item Technical: the selected threshold is persisted in a configuration file and loaded automatically when the system starts.
\end{enumerate}

Acceptance criteria for HU8: Comparative evaluation of pretrained ArcFace and AdaFace
\begin{enumerate}
  \item Functional: ArcFace and AdaFace are run in their pretrained configuration on the LFW, AgeDB, and RFW \textit{benchmarks} and on the degraded variants of LFW and AgeDB, under identical preprocessing and metric conditions, and EER, FAR, FRR, AUC, TAR at fixed FAR, ROC curve, mean inference time, and per-\textit{embedding} memory usage are reported.
  \item Visual: the results are presented as comparison tables and overlaid ROC curves in the documentation and in the \textit{demo}.
  \item Technical: the procedure is reproducible. The version of each model, the inference configuration, the \textit{benchmarks} and their protocol, the random seeds, and the expected metric tolerance are documented. Runs are executed on the reference \textit{hardware} or on a documented equivalent GPU environment.
\end{enumerate}

Acceptance criteria for HU9: Evaluation of pretrained TransFace (optional)
\begin{enumerate}
  \item Functional: if TransFace is included, it is evaluated under identical conditions to ArcFace and AdaFace and reported in an additional row of the comparison tables, with an explicit note that it results from an optional iteration.
  \item Visual: the results are included in the comparison tables and overlaid ROC curves, visually marked as ``optional.''
  \item Technical: the same metrics as in HU8 are reported, and the model version and inference configuration used are documented.
\end{enumerate}

Acceptance criteria for HU10: Preparation of the domain-adaptation \textit{fine-tuning} subset
\begin{enumerate}
  \item Functional: a subset of Glint360K or WebFace4M (5,000--10,000 identities, 100,000--300,000 images) is selected and documented using age-diversity-oriented sampling, and \textit{data augmentation} with synthetic degradations (blur, compression, low resolution) is applied.
  \item Visual: the documentation includes descriptive statistics of the subset (number of identities, age distribution, examples of the degradations applied) and a diagram of the \textit{augmentation} \textit{pipeline}.
  \item Technical: the \textit{train}/\textit{val} split is documented with the identity identifiers assigned to each partition; the overlap between partitions is programmatically verified and is zero; the age distribution per partition is comparable (mean difference $\leq 5$ years).
\end{enumerate}

Acceptance criteria for HU11: Execution of ArcFace's domain-adaptation \textit{fine-tuning}
\begin{enumerate}
  \item Functional: ArcFace is tuned via domain-adaptation \textit{fine-tuning} on the subset prepared in HU10, with a documented scheme for the split, layer \textit{freeze}, \textit{learning rate}, number of epochs, and \textit{batch size}, and intermediate \textit{checkpoints} are preserved.
  \item Visual: the \textit{loss} and validation-metric curves during training are plotted and included in the documentation.
  \item Technical: the effect of \textit{fine-tuning} on the EER for AgeDB and for the degraded variants of LFW and AgeDB, relative to the pretrained model from HU8, is reported as an absolute percentage difference with a 95\% confidence interval estimated via \textit{bootstrap} with $N=1000$ resamples over the evaluation pairs of each \textit{benchmark}; if the effect is not statistically significant, the cause is documented and analyzed (subset size, model capacity, degradation magnitude, \textit{dataset} domain).
\end{enumerate}

Acceptance criteria for HU12: Execution of AdaFace's domain-adaptation \textit{fine-tuning} (optional)
\begin{enumerate}
  \item Functional: if AdaFace tuning is included, it is applied to the same subset prepared in HU10, using the same methodology as HU11.
  \item Visual: the \textit{loss} and validation-metric curves are included in the documentation, labeled as optional.
  \item Technical: the same metrics as in HU11 are reported, using the same statistical methodology (\textit{bootstrap} $N=1000$, 95\% CI).
\end{enumerate}

Acceptance criteria for HU13: Post-\textit{fine-tuning} re-evaluation and comparative analysis
\begin{enumerate}
  \item Functional: the models tuned in HU11 (and, if applicable, HU12) are re-evaluated on LFW, AgeDB, RFW, and the degraded variants under identical conditions to HU8, and tables and figures comparing pretrained vs. tuned are produced.
  \item Visual: the documentation includes overlaid ROC curves, a metrics table by model and \textit{benchmark}, and a chart of the \textit{trade-off} between accuracy (EER) and computational cost (inference time).
  \item Technical: the analysis is stored in a file with the full experimental configuration (model, \textit{benchmark}, threshold, version, seed, split), reproducible from the HU15 \textit{harness}.
\end{enumerate}

Acceptance criteria for HU14: \textit{Liveness} module (\textit{baseline})
\begin{enumerate}
  \item Functional: the system integrates a pretrained \textit{liveness} classifier that, given a facial image, returns a ``genuine face'' probability and applies a threshold to issue the decision.
  \item Visual: the \textit{demo} shows the \textit{liveness} decision alongside the verification decision, and an example attack (printed photo) correctly labeled as an attack in a demonstration run.
  \item Technical: on the selected public \textit{spoofing} \textit{benchmark} (OULU-NPU or CASIA-FASD), the ACER (\textit{Average Classification Error Rate}) or HTER (\textit{Half Total Error Rate}) is reported and compared with the baseline reported in the literature. The choice of thresholds and the \textit{baseline}'s limitations are documented explicitly.
\end{enumerate}

\subsection{Epic 4: Evaluation \textit{harness}, \textit{demo}, and delivery}

Acceptance criteria for HU15: \textit{End-to-end} evaluation \textit{harness}
\begin{enumerate}
  \item Functional: the \textit{harness} accepts the minimum required arguments via the command line (model, \textit{benchmark}, split, seed) and produces an output directory with reproducible metrics (JSON or CSV) and figures (PNG).
  \item Visual: the \textit{harness} documentation (in Markdown) describes the arguments, the expected outputs, and a usage example.
  \item Technical: a full run of the \textit{harness} on the smallest \textit{benchmark} (LFW) completes in a time that is explicitly logged for the reference \textit{hardware} or the GPU environment used (as a rough guide, on the order of tens of minutes on LFW); the actual time is published together with the model version and the inference configuration, so that the verifiable criterion is the documented reproducibility of the run time rather than a fixed, device-dependent threshold. The repository \textit{README} includes the exact commands to reproduce the figures in the report.
\end{enumerate}

Acceptance criteria for HU16: Interactive \textit{demo} and documentation
\begin{enumerate}
  \item Functional: the \textit{demo} in Gradio or Streamlit allows loading an image (or capturing one from the camera), selecting the model, and viewing the verification decision, the similarity level, the \textit{liveness} decision, and, where applicable, the result of the document--face match.
  \item Visual: the \textit{demo} includes a tab or view of experimental results with the ROC curves and the models' comparison tables.
  \item Technical: the \textit{demo} runs locally with a single command, and the documentation describes how to start it. The final project documentation also includes the prototype user manual and the reproduction instructions.
\end{enumerate}

\section{8. CRISP-DM phases}
\label{sec:crispdm}

The project adheres to the \textit{Cross-Industry Standard Process for Data Mining} (CRISP-DM) model, adapted to the specific characteristics of a facial-recognition biometric verification system. The adaptation of each phase to the project context is described below, with emphasis on the experimental phase that forms the core of the work.

\subsection{Phase 1: Business understanding}

Objective: understand the organizational context, the project goals, and the success criteria from a business perspective.

The project is framed by \empclientename's\hspace{1px} need to verify the identity correspondence of its outsourced operational staff. The added value of artificial intelligence lies in automating the facial cross-check, which would otherwise be manual, subjective, and error-prone. The business success metrics are: reduction of identity-impersonation incidents, reduced verification time per shift, and availability of a traceable record of verification attempts. The project's experimental question (\textit{fine-tuning} vs. pretrained, \textit{trade-off} with computational cost) addresses the client's interest in having technical evidence for a possible future deployment decision.

The scope of the prototype is academic: validation will be carried out with public \textit{datasets}, and \textit{fine-tuning} is performed on a subset representative of the domain, not on the company's real data. No operational deployment at the company is planned.

\subsection{Phase 2: Data understanding}

Objective: identify the data sources needed, their nature, and the quality conditions required for \textit{fine-tuning} and evaluation.

The data come exclusively from public facial-recognition \textit{datasets}, selected for their volume, demographic diversity, and license compatibility with academic use. The following will be used:

\begin{itemize}
  \item Glint360K (InsightFace): approximately 17 million images of 360,000 identities, collected from public web sources and distributed for non-commercial research. It will be used as the primary source for \textit{fine-tuning}, on a selected subset.
  \item WebFace4M (Zenodo, 2023): approximately 4 million images of 200,000 identities, available as a modern alternative with a clear license. It will be considered as a complementary source if Glint360K is not available.
  \item LFW (\textit{Labeled Faces in the Wild}): more than 13,000 facial images of 5,749 people, with predefined genuine and adversarial pairs. It will be used as the \textit{baseline} \textit{benchmark} (saturated, expected above 99\%).
  \item AgeDB: a \textit{cross-age} verification \textit{benchmark} with 16,488 images of 568 people, with age ranges up to 101 years. It will be used to assess the models' ability to handle aging between the document photo and the live capture, directly tied to the stated problem.
  \item RFW (\textit{Racial Faces in the Wild}): a \textit{benchmark} specifically designed to measure demographic bias in facial-recognition models. It will be used to report performance by demographic subgroup and to tie the analysis to the data-governance section.
  \item OULU-NPU or CASIA-FASD: public facial-\textit{spoofing} \textit{datasets}, used for evaluating the \textit{liveness} module (\textit{baseline}).
\end{itemize}

Note: \textit{MS-Celeb-1M} (retracted by Microsoft in 2019 due to consent issues) and \textit{CASIA-WebFace} (withdrawn and with an ambiguous license) are explicitly excluded, both of which were mentioned in the original proposal. The licenses of the selected \textit{datasets} will be verified before use and documented in the data-governance section (section 12).

\subsection{Phase 3: Data preparation}

Objective: transform the raw data into inputs suitable for \textit{fine-tuning} and for the \textit{embedding}-generation \textit{pipeline}.

The key transformations are:

\begin{itemize}
  \item Face detection: locate and crop the facial region of each image, removing irrelevant background. Detection models such as MTCNN or RetinaFace will be used.
  \item Facial alignment: correct the orientation of the detected face using \textit{landmark}-based alignment.
  \item Normalization: resize the facial crops to the input resolution expected by each \textit{embedding}-generation model (112$\times$112 pixels for ArcFace and AdaFace).
  \item \textit{Baseline} \textit{embedding} generation: convert each aligned face into a fixed-dimensionality numeric vector (e.g., 512 dimensions) using the pretrained model.
  \item Data splitting: divide the \textit{datasets} into training, validation, and test sets, following the evaluation protocols established for each \textit{benchmark}.
  \item Domain-adaptation \textit{fine-tuning} subset: selection of a subset of Glint360K or WebFace4M (e.g., 5,000--10,000 identities, 100,000--300,000 images) and application of \textit{data augmentation} with synthetic degradations (blur, compression, low resolution) that emulate the target domain's adverse conditions, together with age-diversity-oriented sampling, with a \textit{hold-out} split for validation.
\end{itemize}

\subsection{Phase 4: Modeling}

Objective: select, configure, and tune the facial-\textit{embedding}-generation models.

The problem type is binary identity verification: determining whether two facial images correspond to the same person. This is not identity classification but rather similarity estimation between vector representations.

The algorithms evaluated are:

\begin{itemize}
  \item ArcFace (ResNet-100): a neural network that uses an \textit{Additive Angular Margin Loss} function to generate highly discriminative \textit{embeddings}. The pretrained InsightFace implementation with public \textit{weights} will be used, and domain-adaptation \textit{fine-tuning} will be performed on the subset augmented with synthetic degradations.
  \item AdaFace (CVPR 2022): a neural network that introduces a \textit{quality-adaptive margin} loss function, reporting improvements over ArcFace under adverse conditions. The public pretrained implementation will be used, and it will optionally be tuned on the same subset augmented with synthetic degradations.
  \item TransFace (ICCV 2023, optional): a \textit{transformer}-based architecture for generating facial \textit{embeddings}, reporting incremental improvements over ArcFace on recent \textit{benchmarks}. Its inclusion is conditional on time availability within the schedule and on the success of the main models' domain-adaptation \textit{fine-tuning}. If included, it is evaluated under identical conditions to ArcFace and AdaFace and reported in an additional row of the comparison tables, with an explicit note that it results from an optional iteration and does not affect the work's main findings.
\end{itemize}

The comparison method between \textit{embeddings} is cosine distance. The decision is made via a configurable threshold that balances the false-acceptance rate (FAR) and the false-rejection rate (FRR).

\subsection{Phase 5: Evaluation}

Objective: measure the models' performance and support the choice of the prototype's model, with emphasis on the pretrained-vs-tuned comparison.

The performance metrics are:

\begin{itemize}
  \item \textit{Equal Error Rate} (EER): the point at which FAR and FRR are equal; the goal is to minimize it.
  \item \textit{False Acceptance Rate} (FAR): proportion of adversarial pairs incorrectly accepted.
  \item \textit{False Rejection Rate} (FRR): proportion of genuine pairs incorrectly rejected.
  \item ROC curve (\textit{Receiver Operating Characteristic}): a graphical representation of the trade-off between sensitivity and specificity.
  \item AUC (\textit{Area Under the Curve}): area under the ROC curve; overall discriminative capacity.
  \item TAR at fixed FAR (\textit{True Acceptance Rate} at a given \textit{False Acceptance Rate}, e.g., FAR = $10^{-3}$): genuine acceptance rate at a demanding security operating point; a common reference metric in the biometrics industry.
  \item Mean inference time per \textit{embedding} and memory usage, on the reference \textit{hardware}.
\end{itemize}

The \textit{benchmarks} used are:

\begin{itemize}
  \item LFW: classic \textit{baseline}, expected saturation on both models. It is treated as a saturation test (not a differentiation test between ArcFace and AdaFace), consistent with the redefinition of the experimental design documented in the mitigation for Risk 3 (section 13). Differentiation between models is sought on AgeDB and RFW.
  \item AgeDB: \textit{cross-age}, where most of the \textit{fine-tuning} effect is expected to be observed. It is the primary \textit{benchmark} for experimental success (section 3).
  \item RFW: demographic bias, with per-subgroup reporting.
  \item Degraded variants of LFW and AgeDB: images with synthetically applied blur, compression, and low resolution, emulating the capture-quality heterogeneity of the operational domain. The same degradation operators applied to the \textit{fine-tuning} subset (preparation phase) generate these evaluation variants, so that the training--evaluation loop remains comparable and the improvement (or lack thereof) is interpretable in terms of domain adaptation.
  \item OULU-NPU or CASIA-FASD (for \textit{liveness}): ACER/HTER.
\end{itemize}

\subsection{Phase 6: Deployment}

Objective: integrate the selected model into a demonstrable prototype and document the results reproducibly.

Given the project's academic scope, deployment is limited to:

\begin{itemize}
  \item CLI \textit{harness} or parameterizable notebook that runs the full \textit{pipeline} on the selected \textit{benchmarks} and produces reproducible metrics and figures.
  \item Interactive \textit{demo} with Gradio or Streamlit that showcases the capture and verification flows for the presentation.
  \item Repository with technical documentation, instructions for reproducing the experiments, and evaluation \textit{scripts}.
\end{itemize}

No production deployment, integration with \empclientename's\hspace{1px} real systems, or post-delivery maintenance is planned.

\section{9. Work breakdown structure}
\label{sec:wbs}

Based on each user story defined in section~\ref{sec:backlog}, the work is broken down into concrete, measurable, and time-bounded technical tasks. Each task is estimated in hours and assigned a relative priority (High, Medium, or Low) based on its functional relevance and its impact on the deliverables. Given the individual nature of the project and the granularity of the user stories, technical tasks are estimated in a range of 6 to 30 hours, grouping related subtasks that are executed contiguously. This deviation from the suggested range of 2 to 8 hours is adopted deliberately, in order to preserve traceability between tasks and deliverables without artificially fragmenting the work; non-technical activities, which are continuous in nature, are broken down into blocks associated with each sprint, with estimates of 6 to 32 hours.

The total sum of technical hours (388 h) plus non-technical activities (212 h) amounts to 600 h, which matches the project's total budget distributed across the 12 sprints (section~\ref{sec:sprints}).

\begin{table}[htpb]
  \centering
  \caption{Task breakdown by user story (1/5).}
  \begin{tabularx}{\linewidth}{@{}|X|X|c|c|@{}}
    \hline
    \rowcolor[HTML]{C0C0C0}
    User story & Technical task & Estimate & Priority \\ \hline

    Infra. & Setup of the development environment (Python, PyTorch, CUDA), installation of dependencies (MTCNN, RetinaFace, InsightFace, \textit{adaface-pytorch}), and repository setup & 20 h & High \\ \hline
    HU1 & Implementation of the facial-capture module: camera capture interface, face detection, quality validation (single face, minimum resolution, sharpness), and generation of the reference \textit{embedding} & 28 h & High \\ \cline{2-4}
    HU1 & Implementation of the candidate registration flow: associating the \textit{embedding} with the candidate's data (name, ID number, position) and persistent storage & 10 h & High \\ \hline

    HU2 & Development of the facial-extraction module from a document image: automatic detection of the holder's crop on a national ID, driver's license, or resident card & 14 h & High \\ \cline{2-4}
    HU2 & Validation of the quality of the extracted crop (alignment, lighting, resolution) and generation of the corresponding \textit{embedding} & 8 h & Medium \\ \hline

    HU3 & Implementation of the biometric comparison module: cosine-distance computation between \textit{embeddings}, application of the decision threshold, and issuance of the result (accepted/rejected) & 14 h & High \\ \cline{2-4}
    HU3 & Integration of the full document--face verification flow: live capture, extraction from the document, comparison, and presentation of the result with images and similarity level & 14 h & High \\ \hline

  \end{tabularx}
\end{table}

\begin{table}[htpb]
  \centering
  \caption{Task breakdown by user story (2/5).}
  \begin{tabularx}{\linewidth}{@{}|X|X|c|c|@{}}
    \hline
    \rowcolor[HTML]{C0C0C0}
    User story & Technical task & Estimate & Priority \\ \hline

    HU4 & Implementation of the image-quality validation module: occlusion detection, sharpness estimation (Laplacian variance), lighting assessment, and detected-face count & 18 h & High \\ \cline{2-4}
    HU4 & Design of the error-notification interface: descriptive messages per anomaly type (no face, multiple faces, blurry image, insufficient lighting, partially occluded face) and new-capture request & 14 h & Medium \\ \hline

    HU5 & Implementation of the interactive-time verification \textit{pipeline} at the start of the shift: camera capture, face detection, \textit{embedding} generation, comparison with the stored reference, and issuance of the decision & 16 h & High \\ \cline{2-4}
    HU5 & Development of the shift-flow \textit{demo}: camera preview with detected face, visual accepted/rejected indicator, and similarity level & 12 h & High \\ \hline

    HU6 & Implementation of the retry mechanism: attempt control (maximum of three), logging of each attempt (rejection or invalid), and referral to manual review after the third failure & 10 h & Medium \\ \cline{2-4}
    HU6 & Implementation of the verification-incident log: \textit{timestamp}, worker identifier, failure type, and similarity level for each attempt & 6 h & Medium \\ \hline

  \end{tabularx}
\end{table}

\begin{table}[htpb]
  \centering
  \small
  \renewcommand{\arraystretch}{0.9}
  \caption{Task breakdown by user story (3/5).}
  \begin{tabularx}{\linewidth}{@{}|X|X|c|c|@{}}
    \hline
    \rowcolor[HTML]{C0C0C0}
    User story & Technical task & Estimate & Priority \\ \hline

    HU7 & Design and implementation of the threshold-configuration panel: interface to choose among previously evaluated thresholds and visualization of each threshold's effect (FAR/FRR) & 6 h & High \\ \cline{2-4}
    HU7 & Persistence of the selected threshold in a configuration file and automatic loading when the system starts & 4 h & High \\ \hline

    HU8 & Design of the experimental protocol: definition of LFW, AgeDB, and RFW splits, random seeds, inference configuration (resolution, \textit{batch size}, device), and configuration documentation & 14 h & High \\ \cline{2-4}
    HU8 & Execution of the \textit{baseline} evaluation: runs of pretrained ArcFace and AdaFace on LFW, AgeDB, and RFW and on the degraded variants of LFW and AgeDB, measuring EER, FAR, FRR, AUC, TAR at fixed FAR, ROC curve, inference time, and memory usage & 20 h & High \\ \hline

    HU9 (optional) & Run of pretrained TransFace on the same \textit{benchmarks} to broaden the comparison with a recent \textit{transformer}-based architecture, under identical conditions to ArcFace and AdaFace & 8 h & Low \\ \hline

    HU10 & Preparation of the domain-adaptation \textit{fine-tuning} subset: download, filtering, and \textit{train}/\textit{val} split of a subset of Glint360K (or WebFace4M), application of \textit{data augmentation} with synthetic degradations (blur, compression, low resolution) and age-diversity-oriented sampling, with per-split identity documentation & 8 h & High \\ \hline

    HU11 & Execution of ArcFace's domain-adaptation \textit{fine-tuning}: hyperparameter configuration (\textit{learning rate}, number of epochs, \textit{batch size}, \textit{freeze} scheme), training, and logging of \textit{loss}/metric curves and intermediate \textit{checkpoints} & 30 h & High \\ \hline

    HU12 (optional) & Domain-adaptation \textit{fine-tuning} of AdaFace on the same augmented subset, using the same methodology, if the schedule allows & 12 h & Low \\ \hline

  \end{tabularx}
\end{table}

\begin{table}[htpb]
  \centering
  \small
  \renewcommand{\arraystretch}{0.9}
  \caption{Task breakdown by user story (4/5).}
  \begin{tabularx}{\linewidth}{@{}|X|X|c|c|@{}}
    \hline
    \rowcolor[HTML]{C0C0C0}
    User story & Technical task & Estimate & Priority \\ \hline

    HU13 & Re-evaluation of the tuned models on LFW, AgeDB, and RFW under identical conditions to HU8, storing the results in a reproducible format & 20 h & High \\ \cline{2-4}
    HU13 & Pretrained-vs-tuned comparative analysis: tables, overlaid ROC curves, \textit{trade-off} analysis between EER and inference time, discussion of per-\textit{benchmark} results, and justification of the prototype's model choice & 18 h & High \\ \hline

    HU14 & Integration of the pretrained \textit{liveness} classifier (MiniFASNet or FeatherNet) and evaluation on OULU-NPU (or CASIA-FASD), with ACER/HTER computation and documentation of the \textit{baseline}'s limitations & 18 h & Medium \\ \cline{2-4}
    HU14 & Incorporation of the \textit{liveness} decision into the main \textit{pipeline} and the \textit{demo}, with a labeled attack example (printed photo) & 14 h & Medium \\ \hline

    HU15 & Implementation of the CLI \textit{harness}: argparse with arguments (model, \textit{benchmark}, split, seed), \textit{end-to-end} \textit{pipeline} execution, and saving of metrics and figures to an output directory & 18 h & High \\ \cline{2-4}
    HU15 & \textit{Harness} documentation in Markdown: arguments, expected outputs, usage example, exact commands to reproduce the figures in the report & 6 h & Medium \\ \hline

    HU16 & Implementation of the interactive \textit{demo} with Gradio or Streamlit: tabs for live verification, experimental results, and, where applicable, document--face matching & 12 h & Medium \\ \cline{2-4}
    HU16 & Writing of the prototype user manual and the experiment-reproduction instructions in the repository & 6 h & Medium \\ \hline

    Closeout & Final adjustments to the prototype based on the advisor's feedback & 10 h & Medium \\ \hline

  \end{tabularx}
\end{table}

\begin{table}[htpb]
  \centering
  \caption{Task breakdown by user story (5/5): non-technical activities.}
  \begin{tabularx}{\linewidth}{@{}|X|X|c|c|@{}}
    \hline
    \rowcolor[HTML]{C0C0C0}
    User story & Technical task & Estimate & Priority \\ \hline

    Non-technical & Initial planning: meetings with the advisor and the client, definition and adjustment of the \textit{backlog} and the schedule (Sprint 0) & 30 h & High \\ \hline
    Non-technical & Progress reports and interim reports (Sprints 1 and 5) & 18 h & High \\ \hline
    Non-technical & Partial technical documentation of the \textit{pipeline} (Sprint 3) & 18 h & High \\ \hline
    Non-technical & Follow-up meetings and schedule adjustment (Sprint 4) & 6 h & High \\ \hline
    Non-technical & Training monitoring and logging of \textit{loss} curves and metrics (Sprint 6) & 12 h & High \\ \hline
    Non-technical & Writing of the experimental section of the technical report (Sprint 7) & 12 h & High \\ \hline
    Non-technical & Comprehensive review of the \textit{pipeline} and documentation (Sprint 9) & 26 h & High \\ \hline
    Non-technical & Preparation of the defense presentation and rehearsals (Sprint 10) & 28 h & High \\ \hline
    Non-technical & Final review of the report, delivery of documentation, code, and repository (Sprints 10 and 11) & 32 h & High \\ \hline
    Non-technical & Final meetings, organization of the closing event, and acknowledgments (Sprint 11) & 30 h & High \\ \hline

  \end{tabularx}
\end{table}

The sum of the technical-task estimates amounts to 388 hours and that of the non-technical activities to 212 hours, totaling 600 estimated hours in the breakdown, consistent with the sprint plan and with the budget stated in the project charter. The share of non-technical activities (212 h) exceeds the suggested range of 100 to 120 h due to the weight of the technical report, the bimonthly interim reports, and the defense preparation in an individual project with an experimental core; this is adopted as a justified and deliberate deviation. Additionally, 20 hours of optional tasks are contemplated (evaluation of TransFace as a third evaluation model and \textit{fine-tuning} of AdaFace as a second domain-adaptation pass), present only in the WBS and not in the sprint plan, which will be carried out only if the schedule allows and are not part of the work's mandatory core.

To facilitate traceability between the two breakdowns, the ten non-technical rows in the table above correspond, with the same aggregated hourly estimate per row, to the rows labeled ``Non-tech.'' in the sprint plan (section~\ref{sec:sprints}); each non-technical row indicates in parentheses the sprint in which it is carried out, and its estimate matches the sum of the non-technical rows for the corresponding sprint. Likewise, the technical ``Infra.'' row of the WBS (environment setup, 20 h) corresponds to the ``Planning'' row of Sprint 0, and the technical ``Closeout'' row of the WBS (final adjustments to the prototype, 10 h) corresponds to the ``---'' row of Sprint 10. In this way, the WBS and the sprint plan match exactly at 388 technical hours and 212 non-technical hours (600 h of the mandatory core), differing only by the 20 optional hours referred to in the previous paragraph, which are specific to the WBS.

\section{10. Sprint planning}
\label{sec:sprints}

The project's tasks are organized into 12 sprints (Sprint 0 through Sprint 11) of 2 to 3 weeks each, spread across 4 curricular two-month periods (35 weeks, approximately 8 months, from August 24, 2026 to the end of April 2027). Each sprint concentrates 50 hours of work, resulting in an average load of 20 hours per week (with peaks of 25 hours per week in the 2-week sprints), with transition weeks between periods dedicated to review, writing of interim reports, and meetings with the advisor. The person responsible for all tasks is the author of the final project. The plan includes both technical tasks derived from the user stories and non-technical activities (planning, documentation, meetings, and defense preparation). The experimental core (\textit{fine-tuning} and evaluation on multiple \textit{benchmarks}) is concentrated in sprints 5 through 7, while sprints 1 through 4 correspond to building the base \textit{pipeline} and sprints 8 through 11 to integration, \textit{liveness}, \textit{demo}, documentation, and closeout. The fourth two-month period (sprints 9 through 11) concentrates the \textit{demo}, the comprehensive review, defense preparation, and closeout.

\begin{table}[htpb]
  \centering
  \caption{Sprint plan (1/3).}
  \begin{tabularx}{\linewidth}{@{}|l|X|X|c|c|c|@{}}
    \hline
    \rowcolor[HTML]{C0C0C0}
    Sprint & User story or phase & Task & Hours & Resp. & \% \\ \hline

    Sprint 0 & Non-tech. & Scope, schedule, and \textit{backlog} definition; meetings with the advisor and the client & 30 h & St. & 0 \\ \cline{2-6}
    Sprint 0 & Planning & Environment setup, dependency installation, and repository preparation & 20 h & St. & 0 \\ \hline

    Sprint 1 & HU1 & Implementation of the facial-capture module and generation of the reference \textit{embedding} & 28 h & St. & 0 \\ \cline{2-6}
    Sprint 1 & HU1 & Implementation of the candidate registration flow (data + \textit{embedding} + storage) & 10 h & St. & 0 \\ \cline{2-6}
    Sprint 1 & Non-tech. & Writing of the interim progress report & 12 h & St. & 0 \\ \hline

    Sprint 2 & HU2 & Development of the facial-extraction module from a document image & 14 h & St. & 0 \\ \cline{2-6}
    Sprint 2 & HU2 & Crop-quality validation and generation of the document's \textit{embedding} & 8 h & St. & 0 \\ \cline{2-6}
    Sprint 2 & HU3 & Implementation of the biometric comparison module (cosine distance + threshold) & 14 h & St. & 0 \\ \cline{2-6}
    Sprint 2 & HU3 & Integration of the full document--face verification flow & 14 h & St. & 0 \\ \hline

    Sprint 3 & HU4 & Implementation of the image-quality validation module & 18 h & St. & 0 \\ \cline{2-6}
    Sprint 3 & HU4 & Design of the error-notification interface and new-capture request & 14 h & St. & 0 \\ \cline{2-6}
    Sprint 3 & Non-tech. & Writing of the technical documentation for the registration and verification module & 18 h & St. & 0 \\ \hline

  \end{tabularx}
\end{table}

\begin{table}[htpb]
  \centering
  \caption{Sprint plan (2/3).}
  \begin{tabularx}{\linewidth}{@{}|l|X|X|c|c|c|@{}}
    \hline
    \rowcolor[HTML]{C0C0C0}
    Sprint & User story or phase & Task & Hours & Resp. & \% \\ \hline

    Sprint 4 & HU5 & Implementation of the interactive-time verification \textit{pipeline} at the start of the shift & 16 h & St. & 0 \\ \cline{2-6}
    Sprint 4 & HU5 & Development of the shift-flow \textit{demo} & 12 h & St. & 0 \\ \cline{2-6}
    Sprint 4 & HU6 & Implementation of the retry mechanism and incident logging & 16 h & St. & 0 \\ \cline{2-6}
    Sprint 4 & Non-tech. & Follow-up meetings and schedule adjustment & 6 h & St. & 0 \\ \hline

    Sprint 5 & HU7 & Design and implementation of the threshold-configuration panel & 10 h & St. & 0 \\ \cline{2-6}
    Sprint 5 & HU8 & Design of the experimental protocol and download of \textit{benchmarks} (LFW, AgeDB, RFW) & 14 h & St. & 0 \\ \cline{2-6}
    Sprint 5 & HU8 & Execution of the \textit{baseline} evaluation (pretrained ArcFace, AdaFace) & 20 h & St. & 0 \\ \cline{2-6}
    Sprint 5 & Non-tech. & Writing of the progress report & 6 h & St. & 0 \\ \hline

    Sprint 6 & HU10 & Preparation of the domain-adaptation \textit{fine-tuning} subset (Glint360K or WebFace4M) & 8 h & St. & 0 \\ \cline{2-6}
    Sprint 6 & HU11 & Execution of ArcFace's domain-adaptation \textit{fine-tuning} & 30 h & St. & 0 \\ \cline{2-6}
    Sprint 6 & Non-tech. & Training monitoring and logging of \textit{loss} curves and metrics & 12 h & St. & 0 \\ \hline

  \end{tabularx}
\end{table}

\begin{table}[htpb]
  \centering
  \caption{Sprint plan (3/3).}
  \begin{tabularx}{\linewidth}{@{}|l|X|X|c|c|c|@{}}
    \hline
    \rowcolor[HTML]{C0C0C0}
    Sprint & User story or phase & Task & Hours & Resp. & \% \\ \hline

    Sprint 7 & HU13 & Re-evaluation of tuned models on LFW, AgeDB, and RFW & 20 h & St. & 0 \\ \cline{2-6}
    Sprint 7 & HU13 & Pretrained-vs-tuned comparative analysis, charts, and \textit{trade-off} & 18 h & St. & 0 \\ \cline{2-6}
    Sprint 7 & Non-tech. & Writing of the experimental section of the report & 12 h & St. & 0 \\ \hline

    Sprint 8 & HU14 & Integration of the \textit{liveness} classifier and evaluation on OULU-NPU & 18 h & St. & 0 \\ \cline{2-6}
    Sprint 8 & HU14 & Incorporation of the \textit{liveness} decision into the \textit{pipeline} and the \textit{demo} & 14 h & St. & 0 \\ \cline{2-6}
    Sprint 8 & HU15 & Implementation of the \textit{end-to-end} evaluation CLI \textit{harness} & 18 h & St. & 0 \\ \hline

    Sprint 9 & HU15 & \textit{Harness} documentation and reproduction commands & 6 h & St. & 0 \\ \cline{2-6}
    Sprint 9 & HU16 & Implementation of the interactive \textit{demo} with Gradio/Streamlit & 18 h & St. & 0 \\ \cline{2-6}
    Sprint 9 & Non-tech. & Comprehensive review of the \textit{pipeline} and documentation & 26 h & St. & 0 \\ \hline

    Sprint 10 & Non-tech. & Preparation of the defense presentation & 16 h & St. & 0 \\ \cline{2-6}
    Sprint 10 & Non-tech. & Rehearsals of the presentation with feedback from the advisor & 12 h & St. & 0 \\ \cline{2-6}
    Sprint 10 & --- & Final adjustments to the prototype based on the advisor's feedback & 10 h & St. & 0 \\ \cline{2-6}
    Sprint 10 & Non-tech. & Final review of the report and delivery preparation & 12 h & St. & 0 \\ \hline

    Sprint 11 & Non-tech. & Final meetings with the advisor and the client & 18 h & St. & 0 \\ \cline{2-6}
    Sprint 11 & Non-tech. & Organization of the closing event and acknowledgments & 12 h & St. & 0 \\ \cline{2-6}
    Sprint 11 & Non-tech. & Final delivery of documentation, code, and repository & 20 h & St. & 0 \\ \hline

  \end{tabularx}
\end{table}

The distribution totals 600 hours, spread across 12 sprints (Sprint 0 through Sprint 11) over a period of 35 weeks (4 curricular two-month periods, approximately 8 months), from August 24, 2026 to the end of April 2027. The average load is 20 hours per week (50 h per 2--3-week sprint), with transition weeks between periods. The experimental core (sprints 5 through 7) receives 120 technical hours, reflecting its central role. The fourth two-month period (sprints 9 through 11) concentrates the \textit{demo}, the comprehensive review of the \textit{pipeline} and documentation, defense preparation, and final delivery.

\section{11. Gantt chart (sprints)}
\label{sec:gantt}

Figure~\ref{fig:gantt} presents the project's Gantt chart, based on the sprints defined in section~\ref{sec:sprints}. The horizontal axis represents the project's weeks, spread across 4 curricular two-month periods from August 24, 2026 to the end of April 2027 (35 weeks, approximately 8 months). Each row corresponds to a sprint or a sprint task. Technical tasks derived from user stories are visually distinguished from non-technical activities (planning, documentation, defense). Light gray bands represent transition weeks between periods, dedicated to review, writing of interim reports, and planned rest.

\begin{landscape}
  \begin{figure}[htpb]
    \centering
    \begin{ganttchart}[
        x unit=0.53cm,
        y unit title=0.3cm,
        y unit chart=0.3cm,
        vgrid,
        hgrid,
        title height=2,
        bar/.style={fill=blue!60},
        bar incomplete/.style={fill=blue!30},
        bar label font=\scriptsize,
        title label font=\scriptsize,
        group label font=\scriptsize\bfseries,
        group right shift=0,
        group top shift=0.3,
        group height=.3,
        milestone label font=\scriptsize,
        milestone/.style={fill=red!80}
      ]{1}{35}
      \gantttitle{August}{2}
      \gantttitle{Sept.}{4}
      \gantttitle{Oct.}{4}
      \gantttitle{Nov.}{4}
      \gantttitle{Dec.}{5}
      \gantttitle{Jan.}{4}
      \gantttitle{Feb.}{4}
      \gantttitle{Mar.}{4}
      \gantttitle{Apr.}{4} \\
      \foreach \i in {1,...,35} {
        \gantttitle{\scriptsize S\i}{1}
      } \\

      % Bimestre 1 (semanas 1-8)
      \ganttgroup{Sprint 0: planning}{1}{2} \\
      \ganttbar[bar/.style={fill=orange!60}]{Scope def.}{1}{1} \\
      \ganttbar[bar/.style={fill=orange!60}]{Env. setup}{2}{2} \\

      \ganttgroup{Sprint 1: facial capture}{3}{5} \\
      \ganttbar{HU1: capture + registration}{3}{5} \\
      \ganttbar[bar/.style={fill=orange!60}]{Interim report}{5}{5} \\

      \ganttgroup{Sprint 2: document \& verif.}{6}{8} \\
      \ganttbar{HU2: doc. extraction}{6}{7} \\
      \ganttbar{HU3: biometric comparison}{7}{8} \\

      % Transición bimestre 1-2 (semana 9)
      \ganttbar[bar/.style={fill=gray!20}]{Period 1 transition}{9}{9} \\

      % Bimestre 2 (semanas 10-17)
      \ganttgroup{Sprint 3: quality \& errors}{10}{12} \\
      \ganttbar{HU4: image validation}{10}{11} \\
      \ganttbar{HU4: error interface}{11}{12} \\
      \ganttbar[bar/.style={fill=orange!60}]{Tech. doc.}{12}{12} \\

      \ganttgroup{Sprint 4: shift verif.}{13}{15} \\
      \ganttbar{HU5: verification \textit{pipeline}}{13}{14} \\
      \ganttbar{HU6: retries \& incidents}{15}{15} \\
      \ganttbar[bar/.style={fill=orange!60}]{Follow-up meetings}{15}{15} \\

      \ganttgroup{Sprint 5: \textit{baseline}}{16}{17} \\
      \ganttbar{HU7: threshold config.}{16}{16} \\
      \ganttbar{HU8: \textit{baseline} eval.}{16}{17} \\
      \ganttbar[bar/.style={fill=orange!60}]{Progress report}{17}{17} \\

      % Transición bimestre 2-3 (semana 18)
      \ganttbar[bar/.style={fill=gray!20}]{Period 2 transition}{18}{18} \\

      % Bimestre 3 (semanas 19-26)
      \ganttgroup{Sprint 6: domain adapt. \textit{fine-tuning}}{19}{21} \\
      \ganttbar{HU10: adapt. \textit{fine-tuning} subset}{19}{19} \\
      \ganttbar{HU11: ArcFace \textit{fine-tuning} (domain adapt.)}{20}{21} \\
      \ganttbar[bar/.style={fill=orange!60}]{Training monitoring}{21}{21} \\

      \ganttgroup{Sprint 7: analysis}{22}{24} \\
      \ganttbar{HU13: re-evaluation}{22}{23} \\
      \ganttbar{HU13: comparative analysis}{23}{24} \\
      \ganttbar[bar/.style={fill=orange!60}]{Report (experimental sec.)}{24}{24} \\

      \ganttgroup{Sprint 8: \textit{liveness} + \textit{harness}}{25}{26} \\
      \ganttbar{HU14: \textit{liveness}}{25}{26} \\
      \ganttbar{HU15: CLI \textit{harness}}{26}{26} \\

      % Transición bimestre 3-4 (semana 27)
      \ganttbar[bar/.style={fill=gray!20}]{Period 3 transition}{27}{27} \\

      % Bimestre 4 (semanas 28-34)
      \ganttgroup{Sprint 9: \textit{demo} + review}{28}{30} \\
      \ganttbar{HU15: \textit{harness} doc.}{28}{28} \\
      \ganttbar{HU16: \textit{demo} UI}{28}{30} \\
      \ganttbar[bar/.style={fill=orange!60}]{Comprehensive review}{30}{30} \\

      \ganttgroup{Sprint 10: defense}{31}{32} \\
      \ganttbar[bar/.style={fill=orange!60}]{Presentation prep.}{31}{32} \\
      \ganttbar[bar/.style={fill=orange!60}]{Rehearsals \& adjustments}{32}{32} \\
      \ganttbar[bar/.style={fill=orange!60}]{Report review}{32}{32} \\

      \ganttgroup{Sprint 11: closeout}{33}{34} \\
      \ganttbar[bar/.style={fill=orange!60}]{Final meetings}{33}{33} \\
      \ganttbar[bar/.style={fill=orange!60}]{Final delivery}{34}{34} \\

      \ganttmilestone{Defense}{35}
    \end{ganttchart}
    \caption{Project Gantt chart.}
    \label{fig:gantt}
  \end{figure}
\end{landscape}

Blue represents technical tasks derived from user stories, orange represents non-technical activities (planning, documentation, defense preparation), light gray represents transition weeks between periods dedicated to review, writing of interim reports, and planned rest, and red marks the final-project defense milestone. The sequence reflects the logical dependency between sprints: the capture and extraction tasks (sprints 1--2) are a prerequisite for verification (sprints 2--4), which is in turn a prerequisite for the experimental evaluation (sprints 5--7), the integration of the \textit{liveness} module and the \textit{harness} (sprint 8), the interactive \textit{demo} (sprint 9), and the defense preparation and closeout (sprints 10--11). The public defense is scheduled for \fechaFinalName.

\section{12. Data governance}
\label{sec:gobernanza}

This section provides a comprehensive analysis of the regulatory and ethical framework associated with the use of facial biometric data and the development of an AI-based solution for identity verification.

\subsection{Regulatory compliance}

The project involves facial biometric data (vector \textit{embeddings} derived from facial images), which places it in a particularly sensitive data category from a regulatory standpoint. The applicability of the relevant regulations and the terms of use of the data sources and tools employed are analyzed below.

\subsubsection{Applicable regulations}

Argentine Personal Data Protection Law No. 25,326. Argentina's personal-data-protection law classifies biometric data as sensitive data (art. 2, para. 2), whose processing requires the data subject's express authorization and is subject to additional restrictions. Since the project is carried out within FIUBA (Argentina), this regulation directly applies. However, the project does not collect biometric data from real people: \textit{fine-tuning} and evaluation are carried out exclusively with public \textit{datasets}. Therefore, the consent-collection obligation is not triggered during the experimentation phase. Nonetheless, the prototype will be designed with the law's principles in mind (data quality, purpose limitation, consent, security) as implementation requirements, in order to lay the groundwork for a possible future deployment to comply with the regulation.

General Data Protection Regulation (GDPR). The European GDPR classifies biometric data as a special category of personal data (art. 9), and its processing is prohibited except for certain exceptions (explicit consent, public interest, scientific research). Although the project is carried out in Argentina, mentioning the GDPR is justified for two reasons: (1) \empclientename\hspace{1px} is an international company with operations in Europe, so a possible future deployment of the system could be subject to the GDPR; and (2) some of the public \textit{datasets} used contain images of people from multiple jurisdictions. Since the project operates in an academic context without real data, it does not trigger GDPR obligations, but the prototype's design incorporates principles such as data minimization, the right to erasure, and traceability, which are compatible with the GDPR.

HIPAA (Health Insurance Portability and Accountability Act). This U.S. regulation governs the protection of health data. It is not directly applicable to the project, since the system does not process health information. However, it is mentioned for completeness, since some public facial-recognition \textit{datasets} could contain metadata tied to healthcare contexts at their origin. For the selected \textit{datasets}, this risk does not apply.

\subsubsection{Data sources and licenses}

Glint360K (InsightFace). A \textit{dataset} of approximately 17 million images of 360,000 identities, collected for facial-recognition research by the InsightFace team. It is distributed under an academic-use license. It will be used as the primary source for \textit{fine-tuning}, on a selected subset. The license in effect at the time of download will be verified.

WebFace4M (Zenodo, 2023). A \textit{dataset} of approximately 4 million images of 200,000 identities, publicly available via Zenodo under a clear academic-use license. It will be considered as an alternative to, or complement of, Glint360K if the latter's availability or license terms are not suitable.

LFW (\textit{Labeled Faces in the Wild}). This \textit{dataset} was created by the Computer Vision Laboratory at the University of Massachusetts. The images were obtained from news articles on the web and are distributed for academic, non-commercial research use. No explicit license beyond attribution to the source is required. The use in this project is exclusively academic and evaluative, which is consistent with the \textit{dataset}'s distribution terms.

AgeDB. A \textit{cross-age} verification \textit{dataset} with 16,488 images of 568 people, publicly available for academic use. It will be used as a \textit{benchmark} to assess the models' ability to handle aging.

RFW (\textit{Racial Faces in the Wild}). A \textit{dataset} designed to measure demographic bias in facial-recognition models, publicly available for academic use. It will be used to report performance by demographic subgroup and to tie the analysis to the ethics of biometric-data use.

OULU-NPU or CASIA-FASD. Public facial-\textit{spoofing} \textit{datasets}, widely used in the \textit{liveness}-detection literature. They will be used to evaluate the \textit{liveness} module (\textit{baseline}).

Explicit exclusions. \textit{MS-Celeb-1M} (retracted by Microsoft in 2019 due to consent issues) and \textit{CASIA-WebFace} (withdrawn and with an ambiguous license) are excluded, both of which were mentioned in the original proposal but are not usable in their current state.

\subsubsection{Tools and pretrained models}

ArcFace and AdaFace models. The pretrained implementations are obtained from open-source repositories (\textit{InsightFace}, \textit{adaface-pytorch}). The pretrained weights are distributed under licenses compatible with academic use. ArcFace was published by Imperial College London and its weights are distributed for academic use. AdaFace was published at CVPR 2022 and its weights are distributed under an academic license. No models with restrictive or commercial licenses are used.

Software tools. Development is carried out in Python using open-source libraries: PyTorch for \textit{deep learning}, InsightFace / \textit{adaface-pytorch} for the models, MTCNN and RetinaFace for detection, scikit-learn for evaluation metrics, and Gradio or Streamlit for the interactive \textit{demo}. All the tools are freely usable and compatible with the project's academic context.

\subsubsection{Legal feasibility}

The project is legally feasible for the following reasons: (1) it does not collect biometric data from real people; (2) it uses exclusively public \textit{datasets} under licenses compatible with academic use; (3) the pretrained models are obtained from open-source sources with appropriate licenses; (4) the prototype is not deployed in production nor integrated with the company's real systems. The data-protection measures that would be necessary for a possible future deployment (informed consent, DPIA, deletion and update mechanisms, designation of a data controller) are not part of this prototype and are listed, for informational purposes, in the ``Requirements for a future deployment'' subsection within ``Ethics in the use of artificial intelligence'' (below).

\subsection{Ethics in the use of artificial intelligence}

Biometric identity verification via facial recognition raises significant ethical considerations that must be addressed from the prototype's design stage, even in an academic context.

\subsubsection{Bias in the data and models}

Facial-recognition \textit{datasets} may contain demographic biases that unevenly affect the model's performance across population groups. The biases most relevant to this project are:

\begin{itemize}
  \item Demographic representation bias: public \textit{datasets} may be overrepresented for certain ethnic groups, age ranges, or genders. If the model is evaluated exclusively on these \textit{datasets}, the performance metrics might not reflect actual performance on the diverse population of \empclientename's\hspace{1px} workforce. For example, LFW contains a significantly higher proportion of men than women, and an uneven representation of ethnic groups.
  \item Age bias: \textit{dataset} images may not adequately represent people at extreme age ranges (very young or very old), which could affect the system's accuracy for workforces with a wide age diversity. Evaluation on AgeDB partly addresses this issue by including broad age ranges.
  \item Image-quality bias: \textit{datasets} may be predominantly composed of high-quality images, which does not represent the actual capture conditions in the work environment.
\end{itemize}

These biases do not stem from discriminatory intent but from limitations in the available data sources. Using RFW as one of the project's \textit{benchmarks} makes it possible to quantify the performance gap between demographic subgroups and document it explicitly in the report.

\subsubsection{Social-impact risks}

Misuse of a biometric verification system can create the following risks:

\begin{itemize}
  \item Exclusion or discrimination: if the system shows a significantly higher false-rejection rate for certain demographic groups, the affected workers could face delays, additional reviews, or denial of access, resulting in an uneven work experience.
  \item Use beyond the stated purpose: the verification system could be used for purposes other than identity verification at the start of the shift, such as continuous staff surveillance, movement tracking, or performance evaluation, which would exceed the project's original purpose and could infringe on labor and privacy rights.
  \item Technological dependence: excessive reliance on automated verification could reduce human oversight, creating risks when the system fails or produces ambiguous results.
  \item Improper storage of biometric data: if the facial \textit{embeddings} are not adequately protected, they could be subject to unauthorized access, with serious consequences for workers' privacy, given that biometric data cannot be replaced.
  \item System trivially defeated without robust \textit{liveness}: a facial verification system without \textit{liveness} detection can be fooled with a printed photo of the registered worker, negating its usefulness as a security mechanism.
\end{itemize}

\subsubsection{Mitigation measures}

To address the identified risks, the prototype incorporates the following measures from its design stage, all of which are reflected in the \textit{backlog}, the acceptance criteria, the WBS, or the project's sprints:

\begin{itemize}
  \item Evaluation using standard protocols: the comparative evaluation of models is carried out under established protocols (LFW, AgeDB, RFW) that allow independent reproduction and comparison with the literature. This is realized in HU8 and HU13 and in their associated technical criteria.
  \item Documentation of limitations: the technical report explicitly documents the \textit{datasets}' limitations regarding demographic representation, noting that the performance metrics may not be representative of populations not represented in the evaluation data.
  \item Data minimization: the prototype stores only vector \textit{embeddings} (compact numeric representations), not facial images. The audit logs contain no images or \textit{embeddings}, only verification metadata. This decision is applied in the capture module (HU1) and in the incident log (HU6).
  \item Human review: the system does not issue definitive decisions autonomously. On rejections or ambiguous situations, the flow provides for referral to manual review, keeping the human as the final decision-maker. This decision is realized in HU6 (retry and referral after three rejections) and in HU5 (shift supervisor's decision).
  \item Limited purpose: the system is designed exclusively for identity verification at the start of the shift and during hiring. No continuous real-time recognition, tracking, or crowd-identification functionality is implemented. This limitation is stated explicitly in section 4 (scope).
  \item \textit{Liveness} module as a \textit{baseline}: a pretrained \textit{liveness} classifier is included as a baseline (HU14), with an explicit note that its robust development is left as future work. This approach acknowledges the system's current limitation and opens the door to a later iteration (for example, at a master's-degree stage) that incorporates more modern architectures and domain data.
\end{itemize}

\subsubsection{Requirements for a future deployment}

The following measures are necessary for a possible operational deployment of the system, but they are not part of the scope of this academic prototype and are not implemented in the project's \textit{backlog}, acceptance criteria, WBS, or sprints. They are documented here as explicit guidance for a later iteration, in line with the corresponding exclusion stated in section 4:

\begin{itemize}
  \item Access control for the \textit{embedding} store: authentication and authorization mechanisms restricting queries and modifications to the \textit{embedding} database to authorized users (system administrator, HR team), with access logging.
  \item Record-deletion and update mechanisms: operational implementation of the right to erasure (``right to be forgotten'') and of the rectification or update of the \textit{embeddings} associated with a candidate or worker (e.g., upon reassignment, document change, or withdrawal of consent).
  \item Informed consent of the data subject: a process for collecting, recording, and revoking each worker's express consent for storing their \textit{embedding}, with traceability of acceptances and revocations.
  \item Encryption of \textit{embeddings} at rest and in transit: application of symmetric encryption in the database and TLS in communications between the client and the server, with documented key management.
  \item Formal data-protection procedures: performing a data-protection impact assessment (DPIA), formally designating a data controller, and, where applicable, prior consultation with the supervisory authority.
  \item Production-grade robust \textit{liveness}: replacement of the academic \textit{baseline} (HU14) with a certified anti-\textit{spoofing} system, using training data representative of the domain and evaluation under real conditions.
\end{itemize}

Each of these requirements, when incorporated into the product, should be translated into new user stories, verifiable acceptance criteria, and WBS tasks, maintaining traceability among the purpose, scope, \textit{backlog}, and data governance.

\subsubsection{General ethical reflection}

Biometric identity verification via facial recognition is a technology with significant potential to improve operational security, but it carries inherent risks related to privacy, fairness, and the proportionality of using sensitive data.

In the context of this project, the decision to operate exclusively with public \textit{datasets} and without real worker data is a deliberate ethical choice that allows the system's technical feasibility to be assessed without exposing the privacy of actual individuals. The data-protection measures that would be necessary in a real deployment scenario (informed consent, specific purpose, data minimization, security, and the right to erasure) are explicitly listed in the ``Requirements for a future deployment'' subsection and are excluded from the prototype, in line with section 4.

The project recognizes that automating identity verification should not replace human judgment but complement it. The proposed system issues a recommendation (accepted/rejected) that must be validated by a responsible person, especially in borderline cases. This design decision reflects the principle that artificial intelligence must be subject to human oversight in contexts where decisions directly affect people.

The comparative analysis of models includes the evaluation of demographic biases as part of the selection process, so that the choice of the prototype's model is not based solely on aggregate accuracy metrics but also on the fairness of performance across different population groups.

\section{13. Risk management}
\label{sec:riesgos}

a) Risk identification and estimation of consequences.

Risk 1: insufficient computational capacity on free tiers to run the full \textit{fine-tuning}.

\begin{itemize}
\item Severity (S): 8. Lack of available compute would block delivery of HU11 and, with it, the project's experimental core. It would require reformulating the scope or accepting a partial, unpublishable result.
\item Probability of occurrence (O): 5. Google Colab limits sessions to ~12 h and daily availability; Kaggle offers 30 weekly hours of P100; a serious \textit{fine-tuning} run may require 20-40 h of continuous or accumulated GPU time, with a real risk of falling short or experiencing disconnections.
\end{itemize}

Risk 2: \textit{fine-tuning} does not produce a measurable improvement over the \textit{baseline} on the selected \textit{benchmarks}.

\begin{itemize}
\item Severity (S): 6. If the improvement is null or marginal, the work's experimental core is weakened. It is not blocking, but it requires reorienting the analysis (interpreting why it did not improve, comparison with the literature, ablation).
\item Probability of occurrence (O): 4. Modern models pretrained on large \textit{datasets} already have a high \textit{baseline}; a small subset may not add new information to the model. This is a real risk that must be anticipated and documented.
\end{itemize}

Risk 3: saturated LFW prevents differentiating ArcFace from AdaFace on the \textit{baseline} \textit{benchmark}.

\begin{itemize}
\item Severity (S): 5. Differentiation between the two models would shift toward AgeDB and RFW, which changes the narrative but does not prevent the work. The LFW \textit{benchmark} would remain a saturation test rather than a differentiation test.
\item Probability of occurrence (O): 7. Both ArcFace and AdaFace have performed above 99\% on LFW for years; the probability of obtaining EERs below 0.5\% on both is very high.
\end{itemize}

Risk 4: the main \textit{fine-tuning} \textit{dataset} (Glint360K or WebFace4M) becomes inaccessible or changes license during the project.

\begin{itemize}
\item Severity (S): 7. It blocks HU10 (subset preparation) and HU11 (ArcFace \textit{fine-tuning}) and requires retraining or replacing the subset. The schedule would be affected.
\item Probability of occurrence (O): 2. Both \textit{datasets} are hosted on stable institutional repositories; license changes are unlikely over an 8-month horizon, but not impossible.
\end{itemize}

Risk 5: time exhaustion or \textit{burnout} of the responsible party.

\begin{itemize}
\item Severity (S): 9. The project is individual; if the responsible party cannot dedicate the estimated hours, delivery is delayed or quality is degraded.
\item Probability of occurrence (O): 5. A load of 600 h over 8 months implies an average of 20 h per week with peaks of 25 h in the shorter sprints, ambitious but achievable; personal, work, or academic events can reduce actual availability.
\end{itemize}

Risk 6: the synthetic data for the document--face match (HU2/HU3) is not representative of \empclientename's\hspace{1px} real documents.

\begin{itemize}
\item Severity (S): 5. If the prototype is evaluated with synthetic documents, its usefulness in production would be conditional; however, the project's scope is academic and this is documented explicitly.
\item Probability of occurrence (O): 5. Generating controlled synthetic pairs from public \textit{datasets} is standard practice, but crop quality may degrade for document types not represented.
\end{itemize}

b) Risk-management table (RPN is calculated as RPN = S $\times$ O).

\begin{table}[htpb]
\centering
\begin{tabularx}{\linewidth}{@{}|X|c|c|c|c|c|c|@{}}
\hline
\rowcolor[HTML]{C0C0C0}
Risk & S & O & RPN & S* & O* & RPN* \\ \hline
Insufficient computational capacity on free tiers (R1) & 8 & 5 & 40 & 5 & 2 & 10 \\ \hline
\textit{Fine-tuning} with no measurable improvement over the \textit{baseline} (R2) & 6 & 4 & 24 & 6 & 4 & 24 \\ \hline
Saturated LFW does not differentiate models (R3) & 5 & 7 & 35 & 5 & 7 & 35 \\ \hline
Main \textit{dataset} inaccessible or license change (R4) & 7 & 2 & 14 & 7 & 2 & 14 \\ \hline
Time exhaustion or \textit{burnout} (R5) & 9 & 5 & 45 & 9 & 3 & 27 \\ \hline
Non-representative synthetic data for document--face match (R6) & 5 & 5 & 25 & 5 & 5 & 25 \\ \hline
\end{tabularx}%
\end{table}

c) Adopted criterion.

Mitigation measures will be taken for risks whose original RPN exceeds 30. This applies to R1 and R5. Although R3's original RPN also exceeds 30, it is addressed through a redefinition of the experimental design documented in section 8 (CRISP-DM phases) rather than through a mitigation that reduces its S or O; that is why its revised values S*, O*, and RPN* equal the original ones, but the design decision is justified in the mitigation plan.

Note: values marked with (*) in the table correspond to the estimates after mitigation has been applied. For risks that do not exceed the adopted RPN threshold (R2, R4, and R6), no mitigation plan is drawn up and the original values are repeated.

d) Mitigation plan for the risks that originally exceeded the established maximum RPN.

R1 --- Insufficient computational capacity on free tiers.

\begin{itemize}
\item Action: reduce the size of the \textit{fine-tuning} subset to the minimum that allows a statistically interpretable conclusion (e.g., 50,000--100,000 images instead of 300,000); prioritize \textit{fine-tuning} a single model (ArcFace, HU11) and leave AdaFace's (HU12) as optional; train by layers (freezing the early convolutional layers) to reduce compute; document the time and resources consumed in each run. The auxiliary budget (Colab Pro subscription + cloud \textit{credits}) detailed in table~\ref{tab:presupuesto} financially supports this action.
\item Revised severity (S*): 5. With the reduced subset and the layer-wise scheme, the worst case shifts from ``inability to run the \textit{fine-tuning}'' (S = 8) to ``less conclusive result due to a smaller subset'' (S = 5). Severity drops because \textit{fine-tuning} is no longer impossible to run.
\item Revised occurrence (O*): 2. With the reduced subset, layer \textit{freezing}, and the available auxiliary compute budget (Colab Pro + Kaggle), the required GPU hours fall into a range that runs comfortably; the probability of falling short or experiencing disconnections that block delivery is low. Occurrence drops because the action reduces the probability of running out of compute.
\end{itemize}

R3 --- Saturated LFW does not differentiate models.

\begin{itemize}
\item Action (redefinition of the experimental design): accept LFW as a saturation test (not a differentiation test) and focus the differentiation between ArcFace and AdaFace on AgeDB (\textit{cross-age}) and RFW (demographic bias), where visible gaps are expected. Explicitly report the saturation and discuss its meaning in the report. This decision is documented in section 8 (CRISP-DM phases) as part of the experimental protocol.
\item Revised severity (S*): 5. The action is not a mitigation of the risk in the strict sense (it does not reduce the probability that LFW is saturated nor change its impact on the ability to differentiate models on that \textit{benchmark}); it is a design decision that shifts differentiation to other \textit{benchmarks}. Therefore, S* remains at 5, the same as S.
\item Revised occurrence (O*): 7. For the same reason, the probability that LFW is saturated is unchanged; O* remains at 7, the same as O. RPN* remains at 35.
\end{itemize}

R5 --- Time exhaustion or \textit{burnout}.

\begin{itemize}
\item Action: prioritize the user stories of sprints 5 through 7 (experimental core: HU8, HU10, HU11, HU13) over the user stories of sprints 1 through 4 (base \textit{pipeline}). If the schedule is at risk, reduce the scope of the infrastructure user stories (for example, simplifying HU4's handling of invalid captures) before reducing the experimental scope. Request an extension from the graduate program if the situation requires it.
\item Revised severity (S*): 9. The prioritization reduces the probability of not reaching delivery; it does not change the impact of an eventual delay. Therefore, S* remains at 9, the same as S.
\item Revised occurrence (O*): 3. With clear prioritization of the experimental core and the extension plan as a safety net, the probability of not reaching delivery drops from 5 to 3. Occurrence drops because the action reduces the probability of burnout.
\end{itemize}

\section{14. Sprint Review}
\label{sec:sprint_review}

The \textit{Sprint Review} is a fundamental practice in agile methodologies. It consists of reviewing and assessing what has been completed at the end of a sprint. At this point, progress is presented and it is checked whether the functionality meets the established acceptance criteria. Partial deliverables are also identified and adjustments are considered if needed.

Although the project is still at the planning stage, this section makes it possible to project how the most important \textit{backlog} functionality will be evaluated. This forward-looking view supports value-focused planning and allows reflection on possible obstacles.

Objective: to anticipate how the project's progress will be evaluated as functionality is developed, using at least four user stories from the \textit{product backlog} as a basis.

HU1, HU5, HU8, and HU11 are selected, covering both the base \textit{pipeline} (capture, interactive-time verification) and the experimental core (comparative evaluation and \textit{fine-tuning}).

\begin{table}[htpb]
\small
\renewcommand{\arraystretch}{1.15}
\setlength{\tabcolsep}{4pt}
\begin{tabular}{|>{\raggedright\arraybackslash\hspace{1mm}}p{1.7cm}|
>{\raggedright\arraybackslash}p{2.8cm}|
>{\raggedright\arraybackslash}p{3.1cm}|
>{\raggedright\arraybackslash}p{4.3cm}|
>{\raggedright\arraybackslash}p{2.5cm}|}
\hline
\rowcolor[HTML]{CCCCCC}
\textbf{HU} & \textbf{Associated tasks} & \textbf{Expected deliverable} & \textbf{Concrete evidence of completion} & \textbf{Notes or risks} \\
\hline
HU1 & Facial-capture module and registration flow (38 h, Sprint 1) & Functional capture-and-registration module with stored \textit{embedding} & Test report with 50 captures (10 with one face, 10 with several, 10 blurred, 10 with insufficient lighting, 10 with partial occlusion); \texttt{tests/hu1\_quality.py} with \textit{asserts} on detection confidence $>80\%$ and Laplacian variance $>0.05$; screenshot of a successful registration in the \textit{demo}; \textit{product owner} approval (\textit{issue} closed with the \texttt{approved} label) & Validate with a real camera; camera quality affects the EER of the whole chain \\
\hline
HU5 & Start-of-shift verification \textit{pipeline} (28 h, Sprint 4) & Functional live-verification module with an accepted/rejected indicator & \textit{Log} of 20 runs on the reference \textit{hardware} (2019 16'' MacBook Pro) with total \textit{pipeline} time recorded; \texttt{tests/hu5\_latency.py} verifying time $<5$ s on each run; \textit{screenshot} of the \textit{demo} with the three cases (accepted, rejected, retry); shift supervisor's acceptance in the repository & Capture and processing latency on the reference \textit{hardware} may be adjusted; compare with a GPU run \\
\hline
HU8 & Design of the experimental protocol and \textit{baseline} evaluation of ArcFace and AdaFace on LFW, AgeDB, and RFW (34 h, Sprint 5) & Comparative pretrained tables and ROC curves with EER, FAR, FRR, AUC & File \texttt{results/hu8\_baseline.csv} with per-model, per-\textit{benchmark} metrics; figures \texttt{figures/roc\_*.png} for each \textit{benchmark}; configuration file \texttt{configs/hu8.json} with seeds, splits, and model version; entry in \texttt{CHANGELOG.md} signed off by the advisor & Risk that LFW is saturated and does not differentiate models (R3); rely on AgeDB and RFW \\
\hline
HU11 & \textit{Fine-tuning} of ArcFace on the Glint360K or WebFace4M subset (30 execution hours, Sprint 6) & Tuned ArcFace model with \textit{loss} curves and intermediate \textit{checkpoints} & \path|checkpoints/hu11_arcface_epoch{1..N}.pt| preserved; \path|results/hu11_finetune.csv| with pretrained vs. tuned EER per \textit{benchmark} and 95\% CI (\textit{bootstrap} $N=1000$); figures \path|figures/hu11_loss.png| and \path|figures/hu11_delta_eer.png|; model \textit{release} in the repository tagged \path|hu11-approved| and the advisor's acceptance on the corresponding \textit{issue} & Risk of computational capacity (R1) and of not observing an effect from the tuning (R2) \\
\hline
\end{tabular}
\end{table}

\section{15. Sprint Retrospective}
\label{sec:sprint_retro}

The sprint retrospective is a practice geared toward continuous improvement. At the end of a sprint, the team (or the student, if working individually) reflects on what worked well, what can be improved, and what concrete actions can be implemented to work better in the future.

During the coursework, the use of the Retrospective Star was proposed, which organizes reflection around five axes:

\begin{itemize}
\item What to do more of?
\item What to do less of?
\item What to keep doing?
\item What to start doing?
\item What to stop doing?
\end{itemize}

Even at an early stage, this tool allows the student to plan their way of working, identify potential difficulties ahead of time, and design personal-organization strategies.

Objective: to reflect on the project's initial conditions, identifying strengths, potential difficulties, and improvement strategies, even before development begins.

The table is completed for three technical sprints (Sprint 1 for facial capture, Sprint 6 for \textit{fine-tuning}, and Sprint 9 for the \textit{demo} and review) and one non-technical sprint (Sprint 10 for defense preparation), in line with the project's central milestones.

\begin{table}[htpb]
\small
\renewcommand{\arraystretch}{1.15}
\begin{tabular}{|>{\raggedright\arraybackslash}p{2.2cm}|
>{\raggedright\arraybackslash}p{2.2cm}|
>{\raggedright\arraybackslash}p{2.2cm}|
>{\raggedright\arraybackslash}p{2.2cm}|
>{\raggedright\arraybackslash}p{2.2cm}|
>{\raggedright\arraybackslash}p{2.2cm}|}
\hline
\rowcolor[HTML]{CCCCCC}
\textbf{Sprint type and No.} & \textbf{What to do more of?} & \textbf{What to do less of?} & \textbf{What to keep?} & \textbf{What to start doing?} & \textbf{What to stop doing?} \\
\hline
Technical Sprint 1 (Facial capture) & Testing with a real camera and lighting from the start; documenting design decisions in the repository & Rewriting modules without version control; postponing reading the InsightFace documentation & Small, frequent \textit{commits} with descriptive messages; setting up the environment on a reference machine & Logging \textit{issues} and questions in a project \textit{log}; running tests with synthetic \textit{datasets} before integrating & Changing model or library without assessing the impact \\
\hline
Technical Sprint 6 (\textit{fine-tuning}) & Monitoring \textit{loss} and validation during training; documenting hyperparameters and seeds & Retraining from scratch on every change; experimenting without saving \textit{checkpoints} & Rigorous logging of the experimental configuration; using Weights \& Biases (\textit{wandb}) or an equivalent for \textit{tracking} & Comparing training curves across models; computing the total compute time consumed & Modifying the \textit{learning rate} or \textit{freeze} scheme without documented justification \\
\hline
Technical Sprint 9 (\textit{demo} and review) & End-to-end testing with the \textit{demo}; cross-review of the documentation with the advisor & Postponing known \textit{pipeline} \textit{bugs} to push the \textit{demo} forward; accumulating technical debt without documenting it & Reproducible \textit{demo} with a single command; clear repository \textit{README} & Making short \textit{screencasts} of the \textit{demo} for the defense; asking a third party to run the \textit{harness} from scratch & Debugging live during the defense \\
\hline
Non-technical Sprint 10 (Defense) & Timed rehearsals with an audience; mock committee Q\&A & Rewriting \textit{slides} at the last minute; going into excessive secondary technical detail & Clear structure: problem, solution, experiment, results, limitations; figures with legible labels & Asking a non-specialist to listen to the presentation; preparing answers to frequently asked questions (``why not AdaFace in production?'', ``what effect does \textit{fine-tuning} have on the \textit{baseline}?'') & Adding \textit{slides} or sections that cannot be defended within 15 minutes \\
\hline
\end{tabular}
\end{table}

\end{document}
